\documentclass[11pt,a4paper]{article}

\usepackage[margin=1.05in]{geometry}
\usepackage[T1]{fontenc}
\usepackage[utf8]{inputenc}
\usepackage{lmodern}
\usepackage{microtype}
\usepackage{amsmath,amssymb,amsthm,mathtools,mathrsfs}
\usepackage{enumitem}
\usepackage{xcolor}
\usepackage[colorlinks=true,linkcolor=blue!50!black,citecolor=blue!50!black,urlcolor=blue!50!black]{hyperref}

\numberwithin{equation}{section}
\allowdisplaybreaks

\newtheorem{theorem}{Theorem}[section]
\newtheorem{proposition}[theorem]{Proposition}
\newtheorem{lemma}[theorem]{Lemma}
\newtheorem{corollary}[theorem]{Corollary}
\theoremstyle{definition}
\newtheorem{definition}[theorem]{Definition}
\theoremstyle{remark}
\newtheorem{remark}[theorem]{Remark}

\DeclareMathOperator{\Fix}{Fix}
\DeclareMathOperator{\Ann}{Ann}
\DeclareMathOperator{\Spec}{Spec}
\DeclareMathOperator{\ord}{ord}
\DeclareMathOperator{\Idl}{Idl}
\DeclareMathOperator{\Tail}{Tail}
\DeclareMathOperator{\supp}{supp}
\DeclareMathOperator{\im}{im}
\DeclareMathOperator{\Arg}{Arg}
\DeclareMathOperator{\width}{width}
\DeclareMathOperator{\Rea}{Re}
\DeclareMathOperator{\Ima}{Im}
\DeclareMathOperator{\Res}{Res}

\newcommand{\ZZ}{\mathbb Z}
\newcommand{\NN}{\mathbb N}
\newcommand{\RR}{\mathbb R}
\newcommand{\CC}{\mathbb C}
\newcommand{\PP}{\mathbb P}
\newcommand{\BB}{\mathbf B}
\newcommand{\QQ}{\mathbb Q}
\newcommand{\cO}{\mathcal O}
\newcommand{\cP}{\mathcal P}
\newcommand{\cV}{\mathcal V}
\newcommand{\cA}{\mathcal A}
\newcommand{\cB}{\mathcal B}
\newcommand{\cQ}{\mathcal Q}
\newcommand{\cK}{\mathcal K}
\newcommand{\cI}{\mathcal I}
\newcommand{\cC}{\mathcal C}
\newcommand{\cE}{\mathcal E}
\newcommand{\cS}{\mathcal S}
\newcommand{\taup}{\tau}
\newcommand{\into}{\hookrightarrow}
\newcommand{\onto}{\twoheadrightarrow}
\newcommand{\divv}{\operatorname{div}}
\newcommand{\Sph}{\mathbb S^1}
\newcommand{\ph}{\operatorname{ph}}
\newcommand{\e}{\mathrm e}
\newcommand{\ii}{\mathrm i}

\title{Integrated Split-Zero, Packet, Jet, Phase, and Spectral-Zeta Deconstruction of the Riemann Zeta Function\\
One consolidated note from the divisor-shadow, analytic-completion, packet-geometry, phase-packet, and spectral-regularization strands}
\author{}
\date{May 2026}

\begin{document}
\maketitle

\begin{abstract}
This note consolidates into a single mathematical document the previously separated strands of the uploaded zeta work.  The first strand is the split-zero divisor formalism: the fixed Boolean coefficient, the sheafification of the finite-support presheaf, the comparison with classical Artin thickenings, derivative-enhanced jet stalks, and compactified tail objects at infinity.  The second strand is the full complex-analytic execution for the Riemann zeta function: local expansions at $0$ and $1$, exact trivial-zero derivative formulas, leading jets at non-trivial zeros, centered Hadamard products for the completed zeta function, packet polynomials on the Klein-four orbit space, exact Taylor coefficients, and the packet-grouped partial-fraction expansion of $\zeta'/\zeta$.  The third strand is the Connes--Consani spectral realization and the induced geometry of critical packets and $\zeta$-cycles.  The fourth strand is the Taylor-atomic local phase formalism extracted from Tobley's orthogonal-walk constructions, now applied to the reduced local germ of $\xi$ at a non-trivial zero.  The fifth strand is a corrective appendix to the uploaded modular-Hamiltonian note: the ground-state normalized Wick character, the correct Mellin formula for the twisted spectral Dirichlet series, the failure of the naive identity $\mathcal L_U(0)=-c$, and the precise torus spline/odd-theta/odd-square-zeta bridge.

The integrated output is therefore one note with one notation system and one theorem-lemma-proof chain.  All nontrivial statements included here are either proved in full in the body of the note or are quoted explicitly from standard primary sources when a classical external theorem is used.
\end{abstract}

\tableofcontents

\section{Introduction and integration protocol}

The present document is an integrated replacement for a cluster of earlier notes that had begun to separate into different directions.  The consolidation is organized around five mathematically distinct but compatible packages.

\begin{enumerate}[label=(\arabic*)]
\item The \emph{pointwise split-zero coefficient} of a zeta zero is constant, but the divisor, jet, and compactification structures built from that coefficient are not.  This yields Boolean divisor sheaves, reduced retracts, Artin-thickening comparisons, derivative-enhanced support sheaves, and compactified tails.
\item The \emph{complex-analytic zeta package} refines the purely divisor-theoretic picture.  It supplies the local expansions of $\zeta$, $\Lambda$, and $\xi$, the leading local jets at non-trivial zeros, the exact packet product over the quotient by the Klein-four action generated by $s\mapsto 1-s$ and $s\mapsto \overline s$, the exact Taylor coefficients of the centered completed zeta, and the packet-grouped logarithmic derivative.
\item The \emph{spectral-packet package} places the same non-trivial zeros in the Connes--Consani spectral coordinate and identifies the critical packets with primitive $\zeta$-cycle lengths.
\item The \emph{local phase package} converts the reduced local germ of $\xi$ at a non-trivial zero into a canonical real-analytic phase series, then atomizes that phase into exact normalized affine factors whose finite truncations are orthogonal walks.
\item The \emph{spectral-regularization appendix} corrects the uploaded modular-Hamiltonian note precisely where correction is needed and ties its torus model directly to the odd theta sector and the Dirichlet series $2(1-2^{-2s})\zeta(2s)$.
\end{enumerate}

Two principles govern the integration.

\begin{enumerate}[label=(\alph*)]
\item Whenever two earlier notes treated the same object at different levels of precision, only the more precise formulation is retained as primary, and the other is absorbed as a corollary, refinement, or motivating remark.
\item Whenever an earlier note mixed correct mathematics with an over-strong or incomplete statement, the integrated note keeps the correct part and replaces the defective part by a complete theorem with proof.
\end{enumerate}

The note is self-contained except for three classical external inputs.

\begin{enumerate}[label=(\roman*)]
\item Local analytic geometry on Riemann surfaces, used for the standard identification $\cO_{X,x}\cong \CC\{t\}$ in a centered coordinate \cite{Lee}.
\item The order-one Hadamard factorization of the completed zeta function in centered coordinates, used to pass from the divisor of $\xi$ to a convergent packet product.
\item The Connes--Consani spectral realization of critical zeros and $\zeta$-cycles \cite{ConnesConsaniCycles}.
\end{enumerate}

Everything else is proved directly.

% BEGIN part_refined
\section{Globalization semirings and the fixed Boolean pair}

\begin{definition}
Let $R$ be a commutative ring with identity.  Its \emph{globalization semiring} is
\[
G(R):=R\sqcup\{\taup\},
\]
with operations
\[
a\oplus b:=a+b,\qquad a\oplus\taup=\taup\oplus a:=a,\qquad \taup\oplus\taup:=\taup,
\]
and
\[
a\odot b:=ab,\qquad a\odot\taup=\taup\odot a:=\taup,\qquad \taup\odot\taup:=\taup.
\]
Thus $\taup$ is the additive identity and multiplicative absorber, while $1\in R\subset G(R)$ is the multiplicative identity.
\end{definition}

\begin{definition}
For $u\in R^\times\subset G(R)^\times$, define
\[
\Fix_{G(R)}(u):=\{x\in G(R):u\odot x=x\}.
\]
\end{definition}

\begin{lemma}\label{lem:units-fix}
For every commutative ring $R$,
\[
G(R)^\times=R^\times.
\]
Moreover, for every $u\in R^\times$,
\[
\Fix_{G(R)}(u)=\{\taup\}\cup\Ann_R(u-1).
\]
\end{lemma}

\begin{proof}
If $x\in G(R)$ is invertible, say $x\odot y=1$, then $x\neq\taup$ because $\taup\odot y=\taup\neq1$.  Likewise $y\neq\taup$.  Hence $x,y\in R$ and $xy=1$ inside $R$, so $x\in R^\times$.  The converse is immediate.

Now let $u\in R^\times$.  Since $u\odot\taup=\taup$, the absorber is always fixed.  For $r\in R$ one has
\[
u\odot r=r \iff ur=r \iff (u-1)r=0,
\]
which proves the formula.
\end{proof}

\begin{proposition}\label{prop:split-root-domain}
Let $R$ be an integral domain and let $u\in R^\times$ satisfy $u\neq1$.  Then
\[
\Fix_{G(R)}(u)=\{0,\taup\}.
\]
In particular, if $\xi\in\cO_K^\times$ is a nontrivial root of unity in the ring of integers of a number field $K$, then
\[
\Fix_{G(\cO_K)}(\xi)=\{0,\taup\}.
\]
\end{proposition}

\begin{proof}
By Lemma~\ref{lem:units-fix},
\[
\Fix_{G(R)}(u)=\{\taup\}\cup\Ann_R(u-1).
\]
Because $R$ is a domain and $u-1\neq0$, multiplication by $u-1$ is injective.  Hence $\Ann_R(u-1)=\{0\}$.
\end{proof}

\begin{proposition}\label{prop:boolean-pair}
For every ring $R$, the subset
\[
A_R:=\{0,\taup\}\subset G(R)
\]
is a subsemiring canonically isomorphic to the Boolean semiring $\BB=\{0,1\}$.  Explicitly,
\[
\taup\longleftrightarrow 0\in\BB,
\qquad
0\longleftrightarrow 1\in\BB.
\]
\end{proposition}

\begin{proof}
Inside $A_R$ the operation tables are
\[
0\oplus0=0,\qquad 0\oplus\taup=0,\qquad \taup\oplus\taup=\taup,
\]
and
\[
0\odot0=0,\qquad 0\odot\taup=\taup,\qquad \taup\odot\taup=\taup.
\]
These are exactly the Boolean rules after the displayed identification.
\end{proof}

\begin{corollary}\label{cor:zeta-pointwise}
Let $S=G(\ZZ)$ and $F_1(s)=\zeta(s)-1$.  If $\rho$ is any zero of $\zeta$, then
\[
F_1(\rho)=-1,
\qquad
\Fix_S(F_1(\rho))=\{0,\taup\}\cong\BB.
\]
\end{corollary}

\begin{proof}
The first identity is immediate from $\zeta(\rho)=0$.  The second is Proposition~\ref{prop:split-root-domain} with $R=\ZZ$ and $u=-1$.
\end{proof}

\section{Finite-support presheaves and Boolean divisor shadows}

Throughout this section, $X$ is a connected Riemann surface and
\[
D=\sum_{x\in Z}m_x[x]
\]
is an effective divisor with discrete support $Z\subset X$.  We write $A:=\{0,\taup\}\cong\BB$.

\begin{definition}
The \emph{finite-support presheaf} of split-zero data attached to $D$ is
\[
\cP_D(U):=\bigoplus_{x\in U\cap Z}A^{m_x},
\]
that is, the set of families $(a_x)_{x\in U\cap Z}$ with $a_x\in A^{m_x}$ and all but finitely many $a_x$ equal to the basepoint $(\taup,\dots,\taup)$.  Restriction maps are coordinate projection.

The \emph{Boolean multiplicity shadow} of $D$ is the sheaf of $A$-semimodules
\[
\cV_D(U):=\prod_{x\in U\cap Z}A^{m_x},
\]
and the \emph{reduced Boolean shadow} is
\[
\cA_D(U):=\prod_{x\in U\cap Z}A.
\]
\end{definition}

\begin{lemma}\label{lem:product-sheaf}
Let $Z\subset X$ be closed and discrete, and let $(M_x)_{x\in Z}$ be a family of sets, posets, semimodules, or vector spaces.  Then
\[
\mathcal F_M(U):=\prod_{x\in U\cap Z}M_x
\]
defines a sheaf on $X$.
\end{lemma}

\begin{proof}
Let $U=\bigcup_{i\in I}U_i$ be an open cover.  If two sections of $\mathcal F_M(U)$ agree on each $U_i$, then their $x$-coordinates agree for every $x\in U\cap Z$ because some $U_i$ contains $x$.  Hence the sections are equal.

Conversely, suppose $s_i\in\mathcal F_M(U_i)$ are compatible.  For each $x\in U\cap Z$, choose $i$ with $x\in U_i$ and define the $x$-coordinate of the glued section to be the $x$-coordinate of $s_i$.  Compatibility on overlaps makes this independent of the choice of $i$.  This produces the unique glued section on $U$.
\end{proof}

\begin{theorem}\label{thm:sheafification}
The canonical inclusion of presheaves
\[
\eta_D:\cP_D\longrightarrow \cV_D
\]
exhibits $\cV_D$ as the sheafification of $\cP_D$.
\end{theorem}

\begin{proof}
By Lemma~\ref{lem:product-sheaf}, $\cV_D$ is a sheaf.  It remains to prove the universal property.

Let $\mathcal F$ be any sheaf and let $\phi:\cP_D\to\mathcal F$ be a morphism of presheaves.  We construct a unique sheaf morphism $\widetilde\phi:\cV_D\to\mathcal F$ with $\widetilde\phi\circ\eta_D=\phi$.

Fix an open set $U\subset X$ and a section
\[
s=(s_x)_{x\in U\cap Z}\in\cV_D(U)=\prod_{x\in U\cap Z}A^{m_x}.
\]
For each $x\in U\cap Z$, choose an open neighborhood $U_x\subset U$ with $U_x\cap Z=\{x\}$.  Let $e_x\in\cP_D(U_x)$ be the section whose $x$-coordinate is $s_x$ and whose other coordinates are basepoints.  Set
\[
t_x:=\phi_{U_x}(e_x)\in\mathcal F(U_x).
\]
On an overlap $U_x\cap U_y$ with $x\neq y$, both $e_x$ and $e_y$ restrict to the unique basepoint section because $(U_x\cap U_y)\cap Z=\varnothing$.  Hence $t_x$ and $t_y$ agree on the overlap.  The family $(t_x)_{x\in U\cap Z}$ is therefore compatible, and together with the unique section on $U\setminus Z$ it glues to a section of $\mathcal F(U)$, which we define to be $\widetilde\phi_U(s)$.

This construction is functorial in $U$, so it defines a sheaf morphism $\widetilde\phi$.  By construction, if $s\in\cP_D(U)$ has finite support, the glued section is obtained from the same finitely many local pieces that define $\phi_U(s)$, so $\widetilde\phi_U(\eta_D(s))=\phi_U(s)$.  Thus $\widetilde\phi\circ\eta_D=\phi$.

Finally, uniqueness holds because every section of $\cV_D(U)$ is determined by its restrictions to the open cover
\[
U=(U\setminus Z)\cup\bigcup_{x\in U\cap Z}U_x,
\]
and on each $U_x$ the value of any factorization through a sheaf must agree with $\phi$ on the single-support section carrying the $x$-coordinate.
\end{proof}

\begin{proposition}\label{prop:shadow-stalks}
The sheaves $\cV_D$ and $\cA_D$ are well defined.  Their stalks at $x\in Z$ are
\[
(\cV_D)_x\cong A^{m_x},\qquad (\cA_D)_x\cong A,
\]
and their stalks away from $Z$ are one-point stalks.
\end{proposition}

\begin{proof}
This is Lemma~\ref{lem:product-sheaf}.  Since $Z$ is discrete, for each $x\in Z$ there is an open neighborhood $U$ with $U\cap Z=\{x\}$.  For every smaller neighborhood $V\subset U$ of $x$ one has
\[
\cV_D(V)=A^{m_x},\qquad \cA_D(V)=A,
\]
with identity restrictions, so the direct limits defining the stalks are $A^{m_x}$ and $A$.  If $x\notin Z$, choose $U$ with $U\cap Z=\varnothing$; then all sections on smaller neighborhoods are singleton products.
\end{proof}

\begin{proposition}\label{prop:boolean-cube}
For every integer $m\ge1$, the free rank-$m$ $A$-semimodule $A^m$ is canonically isomorphic to the Boolean cube
\[
B_m:=\cP([m]),\qquad [m]:=\{1,\dots,m\},
\]
by the rule
\[
(a_1,\dots,a_m)\longmapsto \{j:a_j=0\}.
\]
Under this identification, semimodule addition is union of subsets.
\end{proposition}

\begin{proof}
Because $A\cong\BB$, coordinatewise addition on $A^m$ is Boolean OR.  Therefore
\[
\{j:(a_j\oplus b_j)=0\}=\{j:a_j=0\text{ or }b_j=0\},
\]
which is exactly the union of the corresponding subsets.  Bijectivity is immediate.
\end{proof}

\begin{proposition}\label{prop:retract-shadow}
For each $m\ge1$, define
\[
\Sigma_m:A^m\to A,
\qquad
\Sigma_m(a_1,\dots,a_m):=a_1\oplus\cdots\oplus a_m,
\]
and
\[
\Delta_m:A\to A^m,
\qquad
\Delta_m(a):=(a,\dots,a).
\]
Then $\Sigma_m$ and $\Delta_m$ are $A$-semimodule maps and
\[
\Sigma_m\circ\Delta_m=\operatorname{id}_A.
\]
Consequently there are canonical sheaf morphisms
\[
\Sigma_D:\cV_D\to\cA_D,
\qquad
\Delta_D:\cA_D\to\cV_D,
\qquad
\Sigma_D\circ\Delta_D=\operatorname{id}_{\cA_D}.
\]
\end{proposition}

\begin{proof}
If $a=\taup$, then $\Sigma_m(\Delta_m(a))=\taup$.  If $a=0$, then idempotence in $A$ gives $\Sigma_m(\Delta_m(a))=0$.  The global statement is obtained coordinatewise over the support of the divisor.
\end{proof}

\begin{theorem}\label{thm:reduced-vs-multiplicity}
The following are equivalent.
\begin{enumerate}[label=(\roman*)]
\item $\cV_D\cong\cA_D$ in sheaves of $A$-semimodules.
\item $m_x=1$ for every $x\in Z$.
\end{enumerate}
When these conditions fail, $\cA_D$ is nevertheless a canonical split retract of $\cV_D$.
\end{theorem}

\begin{proof}
If every $m_x=1$, then the definitions give $\cV_D=\cA_D$.

Conversely, suppose $\cV_D\cong\cA_D$.  Then the stalks are isomorphic at every $x\in Z$, so
\[
A^{m_x}\cong A
\]
in $A$-semimodules.  Any such isomorphism is in particular a bijection of underlying sets, hence
\[
2^{m_x}=|A^{m_x}|=|A|=2.
\]
Thus $m_x=1$ for all $x$.  The split-retract statement is Proposition~\ref{prop:retract-shadow}.
\end{proof}

\begin{remark}
Because $\Spec(\BB)$ has one point, the Boolean divisor shadow is naturally a sheaf of semiring objects on a zero-dimensional support.  Through the embedding of semirings into blueprints \cite{Lorscheid}, it also admits a blueprint-theoretic interpretation.  We do not use that formalism below.
\end{remark}

\section{Classical Artin thickenings and the chain-cube comparison theorem}

\subsection{Local Artin rings}

\begin{definition}
For $x\in X$ and $m\ge1$, let $\mathfrak m_x\subset\cO_{X,x}$ be the maximal ideal.  Define the local Artin rings
\[
Q_{x,m}:=\cO_{X,x}/\mathfrak m_x^m,
\qquad
K_{x,m}:=\cO_{X,x}/\mathfrak m_x^{m+1}.
\]
For the divisor $D=\sum_{x\in Z}m_x[x]$, set
\[
Q_x:=Q_{x,m_x},\qquad K_x:=K_{x,m_x}.
\]
\end{definition}

\begin{lemma}\label{lem:dvr-local}
Let $x\in X$.  Then there exists a holomorphic coordinate $t$ centered at $x$ such that
\[
\cO_{X,x}\cong\CC\{t\},\qquad \mathfrak m_x=(t).
\]
Consequently
\[
Q_{x,m}\cong\CC\{t\}/(t^m)\cong\CC[t]/(t^m),
\qquad
K_{x,m}\cong\CC\{t\}/(t^{m+1})\cong\CC[t]/(t^{m+1}).
\]
\end{lemma}

\begin{proof}
The existence of a local holomorphic coordinate is part of the definition of a Riemann surface \cite[Chapter~1]{Lee}.  In such a coordinate, the germ ring is the convergent power-series ring $\CC\{t\}$ and the maximal ideal is generated by $t$.  Modulo $(t^m)$ every convergent power series is represented by its degree $<m$ truncation, which gives the polynomial description.
\end{proof}

\begin{proposition}\label{prop:ideal-chain-local}
For every $m\ge1$, the ideals of $Q_{x,m}$ are exactly
\[
I_k:=(t^k)/(t^m),\qquad k=0,1,\dots,m,
\]
where $t$ is any centered local coordinate.  In particular,
\[
Q_{x,m}=I_0\supset I_1\supset\cdots\supset I_m=0
\]
is a chain, independent of the chosen coordinate.
\end{proposition}

\begin{proof}
By Lemma~\ref{lem:dvr-local}, $Q_{x,m}\cong\CC\{t\}/(t^m)$.  Ideals of the quotient correspond to ideals of $\CC\{t\}$ containing $(t^m)$.  Since $\CC\{t\}$ is a discrete valuation ring, every ideal is $(t^k)$ for a unique $k\ge0$.  Those containing $(t^m)$ are precisely $(t^k)$ with $0\le k\le m$, giving the displayed chain.  Because the ideal lattice of a ring is intrinsic, the description is coordinate-independent.
\end{proof}

\begin{definition}
For $m\ge1$, define the terminal-segment chain
\[
C_m:=\{T_r:1\le r\le m+1\}\subset \cP([m]),
\]
where
\[
T_r:=\{r,r+1,\dots,m\}\quad (1\le r\le m),
\qquad
T_{m+1}:=\varnothing.
\]
We also write $B_m:=\cP([m])$ for the full Boolean cube.
\end{definition}

\begin{corollary}\label{cor:ideal-chain-identification}
For every $x\in X$ and $m\ge1$, the ideal lattice $\Idl(Q_{x,m})$ is canonically isomorphic to the chain $C_m$ via
\[
I_k=(t^k)/(t^m)\longmapsto T_{k+1}.
\]
\end{corollary}

\begin{proof}
Proposition~\ref{prop:ideal-chain-local} gives all ideals and their inclusion relations.  The displayed bijection preserves order because
\[
I_k\subseteq I_\ell\iff k\ge\ell
\iff T_{k+1}\subseteq T_{\ell+1}.
\]
\end{proof}

\subsection{The chain-cube comparison}

\begin{definition}
Define the inclusion
\[
i_m:C_m\into B_m
\]
by viewing each terminal segment as a subset of $[m]$.  Define monotone maps
\[
L_m:B_m\to C_m,
\qquad
L_m(I):=\begin{cases}
\varnothing,& I=\varnothing,\\
\{\min I,\min I+1,\dots,m\},& I\neq\varnothing,
\end{cases}
\]
and
\[
R_m:B_m\to C_m,
\qquad
R_m(I):=\bigcup\{T\in C_m:T\subseteq I\}.
\]
\end{definition}

\begin{lemma}\label{lem:LRtails}
For every $I\in B_m$, the sets $L_m(I)$ and $R_m(I)$ lie in $C_m$.  Moreover, $L_m(I)$ is the least terminal segment containing $I$, and $R_m(I)$ is the greatest terminal segment contained in $I$.
\end{lemma}

\begin{proof}
If $I=\varnothing$, then $L_m(I)=\varnothing=T_{m+1}\in C_m$.  If $I\neq\varnothing$, then $L_m(I)=T_{\min I}\in C_m$ and by construction it is the least terminal segment containing $I$.

For $R_m(I)$, note that the union of terminal segments contained in $I$ is again a terminal segment: if $T_r,T_s\subseteq I$, then $T_{\max\{r,s\}}\subseteq I$ and $T_r\cup T_s=T_{\min\{r,s\}}$.  Hence the displayed union belongs to $C_m$ and is plainly the greatest terminal segment contained in $I$.
\end{proof}

\begin{theorem}[Chain-cube comparison]\label{thm:chain-cube}
For every $m\ge1$ one has
\[
L_m\dashv i_m\dashv R_m
\]
in the category of finite posets.  Explicitly, for $I\in B_m$ and $T\in C_m$,
\[
L_m(I)\subseteq T \iff I\subseteq i_m(T),
\]
and
\[
i_m(T)\subseteq I \iff T\subseteq R_m(I).
\]
Hence $C_m$ is both reflective and coreflective in $B_m$.
\end{theorem}

\begin{proof}
By Lemma~\ref{lem:LRtails}, $L_m(I)$ is the least terminal segment containing $I$.  Therefore
\[
L_m(I)\subseteq T\iff I\subseteq T=i_m(T),
\]
which proves $L_m\dashv i_m$.

Similarly, $R_m(I)$ is the greatest terminal segment contained in $I$.  Therefore
\[
i_m(T)=T\subseteq I\iff T\subseteq R_m(I),
\]
which proves $i_m\dashv R_m$.
\end{proof}

\begin{corollary}\label{cor:not-iso-chain-cube}
If $m\ge2$, then $C_m$ and $B_m$ are not isomorphic in finite posets.  Nevertheless they are related by the biadjoint comparison of Theorem~\ref{thm:chain-cube}.
\end{corollary}

\begin{proof}
The chain $C_m$ has width $1$.  The Boolean cube $B_m$ has width at least $2$ because $\{1\}$ and $\{2\}$ are incomparable.  Poset isomorphisms preserve width, so no isomorphism exists when $m\ge2$.
\end{proof}

\subsection{Sheafified comparison}

\begin{definition}
Define sheaves of posets on $X$ by
\[
\cI_D(U):=\prod_{x\in U\cap Z}\Idl(Q_x),
\qquad
\cC_D(U):=\prod_{x\in U\cap Z}C_{m_x},
\qquad
\cB_D(U):=\prod_{x\in U\cap Z}B_{m_x},
\]
with coordinatewise order and restriction maps given by projection.
\end{definition}

\begin{proposition}\label{prop:sheaves-posets}
The assignments $\cI_D$, $\cC_D$, and $\cB_D$ are sheaves of finite posets.  There is a canonical isomorphism
\[
\cI_D\cong\cC_D.
\]
Moreover, under the identification of Proposition~\ref{prop:boolean-cube}, the ordered sheaf underlying the $A$-semimodule sheaf $\cV_D$ is canonically isomorphic to $\cB_D$.
\end{proposition}

\begin{proof}
This is Lemma~\ref{lem:product-sheaf} applied to families of finite posets, together with Corollary~\ref{cor:ideal-chain-identification} and Proposition~\ref{prop:boolean-cube} on each stalk.
\end{proof}

\begin{theorem}\label{thm:sheafified-chain-cube}
There are canonical sheaf morphisms of posets
\[
L_D:\cB_D\to\cC_D,
\qquad
i_D:\cC_D\to\cB_D,
\qquad
R_D:\cB_D\to\cC_D,
\]
obtained coordinatewise from $L_{m_x}$, $i_{m_x}$, and $R_{m_x}$.  Stalkwise one has
\[
(L_D)_x\dashv (i_D)_x\dashv (R_D)_x.
\]
Thus the classical jet-ideal chain sheaf is simultaneously a reflective and a coreflective subobject of the Boolean multiplicity-cube sheaf in sheaves of finite posets.
\end{theorem}

\begin{proof}
Coordinatewise application of Theorem~\ref{thm:chain-cube} commutes with restriction maps because all restrictions are coordinate projections.  Hence the maps are sheaf morphisms.  The stalkwise adjunctions are exactly Theorem~\ref{thm:chain-cube} at multiplicity $m_x$.
\end{proof}

\section{Derivative-enhanced jets and coefficient support}

\subsection{Coordinate-free local jet data}

\begin{definition}
Let $f\in\cO(X)$ be a nonzero holomorphic function with zero divisor
\[
\operatorname{div}(f)=D=\sum_{x\in Z}m_x[x].
\]
For each $x\in Z$, write $[f]_x$ for the class of $f$ in $K_x=\cO_{X,x}/\mathfrak m_x^{m_x+1}$.  Define the sheaves
\[
\cQ_D(U):=\prod_{x\in U\cap Z}Q_x,
\qquad
\cK_D(U):=\prod_{x\in U\cap Z}K_x,
\]
and the distinguished section
\[
[f]_D(U):=([f]_x)_{x\in U\cap Z}\in\cK_D(U).
\]
\end{definition}

\begin{proposition}\label{prop:derivative-enhanced-local}
For each $x\in Z$, the class $[f]_x$ is nonzero and lies in the minimal nonzero ideal
\[
\mathfrak m_x^{m_x}/\mathfrak m_x^{m_x+1}\subset K_x.
\]
Indeed, the principal ideal generated by $[f]_x$ is exactly this minimal nonzero ideal.
\end{proposition}

\begin{proof}
By definition of multiplicity, $\ord_x(f)=m_x$.  Hence in the discrete valuation ring $\cO_{X,x}$ one may write
\[
f=u_x\pi_x^{m_x}
\]
with $u_x$ a unit and $\pi_x$ a uniformizer.  Modulo $\mathfrak m_x^{m_x+1}$ this gives
\[
[f]_x=[u_x]\,[\pi_x]^{m_x},
\]
which is nonzero because $u_x$ remains invertible and $\pi_x^{m_x}\notin\mathfrak m_x^{m_x+1}$.  It lies in $\mathfrak m_x^{m_x}/\mathfrak m_x^{m_x+1}$ by construction.  Since that quotient is one-dimensional over the residue field, every nonzero element generates it, so the ideal generated by $[f]_x$ is exactly the minimal nonzero ideal.
\end{proof}

\begin{corollary}\label{cor:coordinate-derivative}
Let $x\in Z$ and choose a centered local coordinate $t$ at $x$.  Then in $K_x\cong\CC[t]/(t^{m_x+1})$ one has
\[
[f]_x=\lambda_x t^{m_x}
\qquad\text{with}\qquad
\lambda_x\in\CC^\times.
\]
If $X=\CC$ with global coordinate $s$, then
\[
\lambda_x=\frac{f^{(m_x)}(x)}{m_x!}.
\]
\end{corollary}

\begin{proof}
In the chosen coordinate, $f(t)=\lambda_x t^{m_x}+O(t^{m_x+1})$ with $\lambda_x\neq0$ because $\ord_x(f)=m_x$.  Passing to the quotient by $(t^{m_x+1})$ gives the first formula.  On $\CC$, this is Taylor's theorem in the coordinate $s$.
\end{proof}

\subsection{Coefficient-support maps on $\CC$}

From now until the end of this subsection we specialize to $X=\CC$ with global coordinate $s$.

\begin{definition}
For $x\in\CC$ and $m\ge1$, define the derivative-enhanced stalk
\[
E_{x,m}:=(s-x)\CC[s]/(s-x)^{m+1}.
\]
Its distinguished basepoint is the zero class.  Under the basis
\[
(s-x),(s-x)^2,\dots,(s-x)^m,
\]
each element of $E_{x,m}$ can be written uniquely as
\[
\sum_{r=1}^m c_r(s-x)^r.
\]
Define a pointed-set map
\[
\sigma_{x,m}:E_{x,m}\to A^m
\]
by
\[
\sigma_{x,m}\!\left(\sum_{r=1}^m c_r(s-x)^r\right)=(a_1,\dots,a_m),
\qquad
 a_r:=\begin{cases}0,& c_r\neq0,\\ \taup,& c_r=0.\end{cases}
\]
Define also
\[
\iota_{x,m}:A^m\to E_{x,m}
\]
by
\[
\iota_{x,m}(a_1,\dots,a_m):=\sum_{\{r:a_r=0\}}(s-x)^r.
\]
\end{definition}

\begin{proposition}\label{prop:pointed-retract}
The maps $\sigma_{x,m}$ and $\iota_{x,m}$ are pointed-set morphisms and satisfy
\[
\sigma_{x,m}\circ\iota_{x,m}=\operatorname{id}_{A^m}.
\]
Consequently, if
\[
\cE_D(U):=\prod_{x\in U\cap Z}E_{x,m_x}
\]
on $\CC$, then the coordinatewise maps
\[
\sigma_D:\cE_D\to\cV_D,
\qquad
\iota_D:\cV_D\to\cE_D
\]
exhibit the underlying pointed-set sheaf of $\cV_D$ as a retract of $\cE_D$.
\end{proposition}

\begin{proof}
Both maps preserve the chosen basepoints by construction.  If $a=(a_1,\dots,a_m)\in A^m$, then the coefficient of $(s-x)^r$ in $\iota_{x,m}(a)$ is nonzero exactly when $a_r=0$.  Applying $\sigma_{x,m}$ therefore recovers $a$.  The sheaf statement is obtained coordinatewise and uses Lemma~\ref{lem:product-sheaf}.
\end{proof}

\begin{proposition}\label{prop:not-cmon-iso}
If $m\ge1$, then $E_{x,m}$ and $A^m$ are not isomorphic as commutative monoids under addition.
\end{proposition}

\begin{proof}
The additive monoid of $A^m$ is idempotent because addition is coordinatewise Boolean OR:
\[
u+u=u\qquad (u\in A^m).
\]
The additive monoid of $E_{x,m}$ is not idempotent: for example,
\[
(s-x)+(s-x)=2(s-x)\neq (s-x)
\]
in characteristic $0$.  Idempotence of the entire monoid is preserved by monoid isomorphism, so no isomorphism exists.
\end{proof}

\begin{corollary}\label{cor:zeta-support-slot}
Let $\rho$ be a zero of zeta of multiplicity $m_\rho$.  In the stalk
\[
E_\rho:=E_{\rho,m_\rho}
\]
one has
\[
[\zeta]_\rho=\lambda_\rho (s-\rho)^{m_\rho},
\qquad
\lambda_\rho:=\frac{\zeta^{(m_\rho)}(\rho)}{m_\rho!}\neq0.
\]
Under the support map $\sigma_{\rho,m_\rho}$, this class is sent to the element of $A^{m_\rho}$ whose only nonbasepoint coordinate is the terminal one.  Equivalently, it picks the atom $\{m_\rho\}\subset [m_\rho]$ of the terminal-segment chain.
\end{corollary}

\begin{proof}
The displayed formula is Corollary~\ref{cor:coordinate-derivative} applied to $f=\zeta$.  Since only the coefficient of $(s-\rho)^{m_\rho}$ is nonzero, the support vector has a single nonbasepoint entry in the terminal slot.
\end{proof}

\section{Compactification and tail objects at infinity}

Let
\[
j:\CC\into\PP^1(\CC)
\]
be the standard open immersion.

\subsection{Pointed-set tails}

\begin{definition}
Let $Z\subset\CC$ be closed and discrete, and let $(M_x)_{x\in Z}$ be pointed sets with basepoints $0_x$.  Define the sheaf
\[
\mathcal F_M(U):=\prod_{x\in U\cap Z}M_x
\qquad (U\subset\CC\text{ open}).
\]
For an open set $U\subset\PP^1(\CC)$ containing $\infty$, write $U^\circ:=U\cap\CC$.  Define
\[
\Tail((M_x)_{x\in Z}):=\prod_{x\in Z}M_x\Big/\sim,
\]
where
\[
(a_x)\sim(b_x)\iff a_x=b_x\text{ for all but finitely many }x\in Z.
\]
\end{definition}

\begin{theorem}\label{thm:stalks-at-infty}
With the notation above, the canonical morphism
\[
j_!\mathcal F_M\longrightarrow j_*\mathcal F_M
\]
has the following stalks.
\begin{enumerate}[label=(\roman*)]
\item If $x\in Z$, then
\[
(j_!\mathcal F_M)_x\cong (j_*\mathcal F_M)_x\cong M_x.
\]
\item If $y\in\CC\setminus Z$, then
\[
(j_!\mathcal F_M)_y\cong (j_*\mathcal F_M)_y\cong *.
\]
\item At $\infty$ one has
\[
(j_!\mathcal F_M)_\infty\cong *,
\qquad
(j_*\mathcal F_M)_\infty\cong \Tail((M_x)_{x\in Z}).
\]
\end{enumerate}
Consequently, if the tail pointed set is nontrivial, then
\[
j_!\mathcal F_M\not\cong j_*\mathcal F_M
\]
in sheaves of pointed sets on $\PP^1(\CC)$.
\end{theorem}

\begin{proof}
The first two assertions are immediate from sufficiently small neighborhoods in $\CC$: around a point of $Z$ the support is a singleton, while around a point of $\CC\setminus Z$ the support is empty.

For the stalk of $j_!\mathcal F_M$ at $\infty$, note that a section of $j_!\mathcal F_M(U)$ on a neighborhood $U$ of $\infty$ is, by definition, eventually equal to the basepoint outside a finite subset of $U^\circ\cap Z$.  Shrinking $U$ so that those finitely many points disappear shows that every germ becomes the basepoint germ.

For the stalk of $j_*\mathcal F_M$ at $\infty$, a germ is represented by a section on some neighborhood $U$ of $\infty$, hence by a family on the cofinite subset $U^\circ\cap Z\subset Z$.  Extend that family by basepoints on the omitted finite subset of $Z$.  This identifies the germ with a class in the quotient modulo finite modification.  Two such families define the same germ exactly when they agree on a smaller neighborhood of $\infty$, equivalently when they differ only at finitely many indices.  This gives the stated identification.

If the stalks at $\infty$ are not isomorphic, the sheaves are not isomorphic.
\end{proof}

\subsection{Vector-space tails}

\begin{definition}
Let $(V_x)_{x\in Z}$ be a family of nonzero complex vector spaces.  Define the sheaf
\[
\mathcal F_V(U):=\prod_{x\in U\cap Z}V_x
\qquad (U\subset\CC\text{ open})
\]
and the tail vector space
\[
\Tail_{\mathrm{vec}}((V_x)_{x\in Z}):=\prod_{x\in Z}V_x\Big/\bigoplus_{x\in Z}V_x.
\]
\end{definition}

\begin{proposition}\label{prop:vector-tail-exact}
There is a short exact sequence of sheaves of complex vector spaces on $\PP^1(\CC)$,
\[
0\longrightarrow j_!\mathcal F_V
\longrightarrow j_*\mathcal F_V
\longrightarrow i_{\infty *}\Tail_{\mathrm{vec}}((V_x)_{x\in Z})
\longrightarrow0.
\]
If $Z$ is infinite, then the tail vector space is nonzero.  Moreover,
\[
H^q\bigl(\PP^1(\CC),j_!\mathcal F_V\bigr)=0
\qquad (q\ge1).
\]
\end{proposition}

\begin{proof}
On an open set $U\subset\PP^1(\CC)$ not containing $\infty$, the sheaves $j_!\mathcal F_V$ and $j_*\mathcal F_V$ coincide.  If $\infty\in U$, then
\[
(j_*\mathcal F_V)(U)=\prod_{x\in U^\circ\cap Z}V_x,
\qquad
(j_!\mathcal F_V)(U)=\bigoplus_{x\in U^\circ\cap Z}V_x,
\]
because sections of $j_!$ are precisely those with finite support near $\infty$.  Since $U^\circ\cap Z$ is cofinite in $Z$, the quotient is canonically identified with
\[
\prod_{x\in Z}V_x\Big/\bigoplus_{x\in Z}V_x.
\]
This proves exactness and identifies the quotient with a skyscraper at $\infty$.

If $Z$ is infinite and each $V_x\neq0$, any family $(v_x)$ with all $v_x\neq0$ defines a nonzero class in the quotient, so the tail vector space is nonzero.

The sheaves $j_*\mathcal F_V$ and $i_{\infty*}\Tail_{\mathrm{vec}}((V_x))$ are flabby: restriction maps on products are coordinate projections, hence surjective, and the same is immediate for a skyscraper sheaf.  The long exact sequence in sheaf cohomology therefore gives
\[
H^q(\PP^1(\CC),j_!\mathcal F_V)=0\qquad(q\ge2),
\]
and an exact segment
\[
0\to \Gamma(j_!\mathcal F_V)\to \Gamma(j_*\mathcal F_V)\to \Tail_{\mathrm{vec}}((V_x))\to H^1(\PP^1(\CC),j_!\mathcal F_V)\to0.
\]
But
\[
\Gamma(j_!\mathcal F_V)=\bigoplus_{x\in Z}V_x,
\qquad
\Gamma(j_*\mathcal F_V)=\prod_{x\in Z}V_x,
\]
and the middle map is the quotient map defining the tail vector space, hence is surjective.  Therefore $H^1$ vanishes as well.
\end{proof}

\section{The zeta divisor, packet descent, and trivial-zero weights}

\subsection{The zeta divisor}

\begin{definition}
Let
\[
D_\zeta:=\sum_{\rho:\,\zeta(\rho)=0}m_\rho[\rho]
\]
be the effective divisor of zeta zeros on $\CC$.  Set
\[
\cP_\zeta:=\cP_{D_\zeta},\quad
\cV_\zeta:=\cV_{D_\zeta},\quad
\cA_\zeta:=\cA_{D_\zeta},\quad
\cI_\zeta:=\cI_{D_\zeta},\quad
\cQ_\zeta:=\cQ_{D_\zeta},\quad
\cK_\zeta:=\cK_{D_\zeta}.
\]
\end{definition}

\begin{corollary}\label{cor:zeta-shadow-sheaves}
At every zero $\rho$ of zeta,
\[
(\cV_\zeta)_\rho\cong A^{m_\rho},
\qquad
(\cA_\zeta)_\rho\cong A,
\qquad
\Fix_S(F_1(\rho))=A.
\]
The reduced zeta shadow is a canonical split retract of the multiplicity shadow.
\end{corollary}

\begin{proof}
Combine Proposition~\ref{prop:shadow-stalks}, Proposition~\ref{prop:retract-shadow}, and Corollary~\ref{cor:zeta-pointwise}.
\end{proof}

\subsection{Completed zeta and packet descent}

\begin{definition}
Define the completed zeta function
\[
\xi(s):=\frac12 s(s-1)\pi^{-s/2}\Gamma\!\left(\frac s2\right)\zeta(s).
\]
Let $D_\xi$ be its zero divisor.  Thus $D_\xi$ is the divisor of the nontrivial zeros of zeta.
\end{definition}

\begin{definition}
Let
\[
c(s):=\overline s,
\qquad
w(s):=1-s,
\qquad
W:=\{\mathrm{id},c,w,cw\}\cong C_2\times C_2.
\]
\end{definition}

\begin{lemma}\label{lem:xi-symmetry-multiplicity}
If $\rho$ is a zero of $\xi$, then $\overline\rho$, $1-\rho$, and $1-\overline\rho$ are also zeros of $\xi$, and all four points occur with the same multiplicity.
\end{lemma}

\begin{proof}
The functional equation $\xi(1-s)=\xi(s)$ gives the reflection symmetry.  The Dirichlet series for $\zeta$ has real coefficients on $\Re(s)>1$, so $\overline{\zeta(\overline s)}=\zeta(s)$ there, and analytic continuation yields the same relation meromorphically on all of $\CC$.  The elementary prefactor in $\xi$ is also conjugation-compatible, so $\overline{\xi(\overline s)}=\xi(s)$.  Therefore zeros are preserved by conjugation.  Multiplicities are preserved because both symmetries are local analytic automorphisms with nonzero derivative in absolute value.
\end{proof}

\begin{definition}
Let $Y_\xi:=|D_\xi|/W$ be the orbit space of nontrivial-zero packets, equipped with the quotient topology.  Since $|D_\xi|$ is discrete, so is $Y_\xi$.  For an orbit $O\in Y_\xi$, write $m_O$ for the common multiplicity of its points.

Define the \emph{Boolean packet sheaf}
\[
\mathcal P_\xi^{\cV}(U):=\prod_{O\in U}A^{m_O}
\]
and the \emph{chain packet sheaf}
\[
\mathcal P_\xi^{\cC}(U):=\prod_{O\in U}C_{m_O}
\]
on open sets $U\subset Y_\xi$.
\end{definition}

\begin{theorem}\label{thm:packet-descent}
Let $q:|D_\xi|\to Y_\xi$ be the quotient map.  The group $W$ acts naturally on the sheaves $\cV_{D_\xi}$ and $\cC_{D_\xi}$.  For each orbit $O\in Y_\xi$ of cardinality $r_O$, the stalk of $q_*\cV_{D_\xi}$ at $O$ is canonically
\[
(q_*\cV_{D_\xi})_O\cong (A^{m_O})^{r_O},
\]
and its $W$-fixed subsemimodule is the diagonal copy of $A^{m_O}$.  Hence
\[
(q_*\cV_{D_\xi})^W\cong \mathcal P_\xi^{\cV}.
\]
Likewise,
\[
(q_*\cC_{D_\xi})_O\cong (C_{m_O})^{r_O},
\qquad
(q_*\cC_{D_\xi})^W\cong \mathcal P_\xi^{\cC}.
\]
Thus both the Boolean multiplicity shadow and the classical jet-ideal chain descend to packet sheaves on the orbit space.
\end{theorem}

\begin{proof}
Because $Y_\xi$ is discrete, the stalk of $q_*\cV_{D_\xi}$ at an orbit $O$ is the product of the stalks over the points in that orbit.  Lemma~\ref{lem:xi-symmetry-multiplicity} gives the common multiplicity $m_O$, so
\[
(q_*\cV_{D_\xi})_O\cong \prod_{\rho\in O}A^{m_O}\cong (A^{m_O})^{r_O}.
\]
The group $W$ acts by permuting the $r_O$ factors transitively.  A tuple in $(A^{m_O})^{r_O}$ is invariant under this action if and only if all coordinates are equal.  Thus the fixed part is the diagonal copy of $A^{m_O}$.  Since the orbit space is discrete, these diagonal stalks assemble into the sheaf $\mathcal P_\xi^{\cV}$.

The chain case is identical, replacing $A^{m_O}$ by $C_{m_O}$.
\end{proof}

\begin{remark}
No hypothesis such as RH is used here.  If RH holds, every nontrivial-zero packet has cardinality $2$; off the critical line, the packet size is $4$.
\end{remark}

\subsection{Trivial zeros}

\begin{proposition}\label{prop:trivial-simple}
For every integer $n\ge1$, the trivial zero $-2n$ of zeta is simple.
\end{proposition}

\begin{proof}
The functional equation reads
\[
\zeta(s)=2^s\pi^{s-1}\sin\!\left(\frac{\pi s}{2}\right)\Gamma(1-s)\zeta(1-s).
\]
At $s=-2n$, the factors $2^s$, $\pi^{s-1}$, and $\Gamma(1-s)=\Gamma(1+2n)$ are holomorphic and nonzero, while $\zeta(1-s)=\zeta(2n+1)\neq0$.  The only vanishing factor is the sine term, and it has a simple zero because
\[
\frac{d}{ds}\sin\!\left(\frac{\pi s}{2}\right)\Big|_{s=-2n}=\frac\pi2\cos(-n\pi)=\frac\pi2(-1)^n\neq0.
\]
\end{proof}

\begin{corollary}\label{cor:trivial-stalk}
On the trivial-zero sector, the Boolean multiplicity stalk is just $A$:
\[
(\cV_\zeta)_{-2n}\cong A\qquad(n\ge1).
\]
\end{corollary}

\begin{proof}
This is Proposition~\ref{prop:shadow-stalks} together with Proposition~\ref{prop:trivial-simple}.
\end{proof}

\begin{proposition}\label{prop:trivial-derivative}
For every integer $n\ge1$,
\[
\zeta'(-2n)=(-1)^n\frac{(2n)!}{2(2\pi)^{2n}}\,\zeta(2n+1).
\]
Hence
\[
|\zeta'(-2n)|=\frac{(2n)!}{2(2\pi)^{2n}}\,\zeta(2n+1).
\]
\end{proposition}

\begin{proof}
Set
\[
g_n(s):=2^s\pi^{s-1}\Gamma(1-s)\zeta(1-s).
\]
Then $\zeta(s)=g_n(s)\sin(\pi s/2)$ and $g_n$ is holomorphic at $s=-2n$ with
\[
g_n(-2n)=2^{-2n}\pi^{-2n-1}(2n)!\,\zeta(2n+1).
\]
Differentiating and using that the sine factor vanishes at $-2n$ gives
\[
\zeta'(-2n)=g_n(-2n)\cdot \frac\pi2\cos(-n\pi)
=2^{-2n}\pi^{-2n-1}(2n)!\,\zeta(2n+1)\cdot\frac\pi2(-1)^n,
\]
which simplifies to the stated formula.  Since $\zeta(2n+1)>0$, the absolute-value formula follows.
\end{proof}

\begin{definition}
For each trivial zero $-2n$, define its \emph{canonical Boolean weight}
\[
\omega_{-2n}:A\to\RR_{\ge0}
\]
by
\[
\omega_{-2n}(\taup):=0,
\qquad
\omega_{-2n}(0):=|\zeta'(-2n)|.
\]
\end{definition}

\begin{corollary}\label{cor:trivial-weight}
The unique nonbasepoint of the trivial Boolean stalk carries the exact analytic weight
\[
\omega_{-2n}(0)=\frac{(2n)!}{2(2\pi)^{2n}}\,\zeta(2n+1).
\]
\end{corollary}

\begin{proof}
This is immediate from the definition and Proposition~\ref{prop:trivial-derivative}.
\end{proof}

\subsection{Compactified zeta tails}

\begin{definition}
Define the zeta tail objects
\[
T_\zeta^{\cV}:=\Tail\bigl((A^{m_\rho})_{\rho:\,\zeta(\rho)=0}\bigr),
\qquad
T_\zeta^{\cA}:=\Tail\bigl((A)_{\rho:\,\zeta(\rho)=0}\bigr),
\]
and the jet-vector-space tails
\[
T_\zeta^{\cQ}:=\Tail_{\mathrm{vec}}\bigl((Q_\rho)_{\rho:\,\zeta(\rho)=0}\bigr),
\qquad
T_\zeta^{\cK}:=\Tail_{\mathrm{vec}}\bigl((K_\rho)_{\rho:\,\zeta(\rho)=0}\bigr).
\]
\end{definition}

\begin{theorem}\label{thm:zeta-tail}
The compactified Boolean and jet shadows of zeta satisfy:
\begin{enumerate}[label=(\roman*)]
\item the stalks at $\infty$ are
\[
(j_*\cV_\zeta)_\infty\cong T_\zeta^{\cV},
\qquad
(j_*\cA_\zeta)_\infty\cong T_\zeta^{\cA},
\]
\[
(j_*\cQ_\zeta)_\infty\cong T_\zeta^{\cQ},
\qquad
(j_*\cK_\zeta)_\infty\cong T_\zeta^{\cK};
\]
\item all four tail objects are nonzero;
\item for the jet-vector-space shadows there are short exact sequences
\[
0\to j_!\cQ_\zeta\to j_*\cQ_\zeta\to i_{\infty*}T_\zeta^{\cQ}\to0,
\]
\[
0\to j_!\cK_\zeta\to j_*\cK_\zeta\to i_{\infty*}T_\zeta^{\cK}\to0,
\]
and therefore
\[
H^q\bigl(\PP^1(\CC),j_!\cQ_\zeta\bigr)=H^q\bigl(\PP^1(\CC),j_!\cK_\zeta\bigr)=0
\qquad(q\ge1).
\]
\end{enumerate}
\end{theorem}

\begin{proof}
Part~(i) is Theorem~\ref{thm:stalks-at-infty} for Boolean shadows and Proposition~\ref{prop:vector-tail-exact} for jet shadows.  Since zeta has infinitely many trivial zeros, the index set is infinite; the class of the everywhere-nonbasepoint family on the trivial-zero sector shows that the Boolean tails are nonzero, and any everywhere-nonzero vector in the product gives nonzero jet tails.  This proves~(ii).  Part~(iii) is Proposition~\ref{prop:vector-tail-exact}.
\end{proof}
% END part_refined

\section{Classical functional equation and completed functions}

To pass from the divisor-shadow formalism to the full analytic structure of the Riemann zeta function we fix the standard completed functions
\[
\Lambda(s):=\pi^{-s/2}\Gamma\!\left(\frac{s}{2}\right)\zeta(s),
\qquad
\xi(s):=\frac12 s(s-1)\Lambda(s).
\]
The functional equation of $\zeta$ will be used in the form
\begin{equation}\label{eq:functional-zeta}
\zeta(s)=2^s\pi^{s-1}\sin\!\left(\frac{\pi s}{2}\right)\Gamma(1-s)\zeta(1-s).
\end{equation}
The trivial zeros occur at the negative even integers, the non-trivial zeros are exactly the zeros of $\xi$, and the action of the Klein four-group
\[
W=\{\mathrm{id},\ s\mapsto 1-s,\ s\mapsto \overline s,\ s\mapsto 1-\overline s\}
\]
organizes the non-trivial zeros into packets on the quotient space $Y_\xi$ introduced above.

% BEGIN part_analytic
\subsection{The pole at $s=1$ and the regular point $s=0$}

\begin{proposition}\label{prop:stieltjes-expansion}
There is a unique sequence $(\gamma_n)_{n\ge 0}$ of complex numbers such that, in a punctured neighborhood of $s=1$,
\begin{equation}\label{eq:stieltjes-expansion}
\zeta(s)=\frac{1}{s-1}+\sum_{n=0}^{\infty}\frac{(-1)^n\gamma_n}{n!}(s-1)^n.
\end{equation}
Equivalently,
\[
\zeta(s)-\frac{1}{s-1}
\]
is holomorphic at $s=1$ and its Taylor coefficients are the Stieltjes constants.
\end{proposition}

\begin{proof}
Since $\zeta$ has a simple pole at $s=1$ with residue $1$, the function
\[
h(s):=\zeta(s)-\frac{1}{s-1}
\]
is holomorphic near $s=1$.  Therefore it has a unique Taylor expansion
\[
h(s)=\sum_{n=0}^{\infty} c_n(s-1)^n.
\]
Writing $c_n=(-1)^n\gamma_n/n!$ gives the asserted form.
\end{proof}

\begin{proposition}\label{prop:zeta-at-zero}
Near $s=0$ one has the exact expansion
\begin{equation}\label{eq:zeta-zero-expansion}
\zeta(s)=-\frac12-\frac12\log(2\pi)s+O(s^2).
\end{equation}
Consequently,
\[
\zeta(0)=-\frac12,
\qquad
\zeta'(0)=-\frac12\log(2\pi).
\]
Moreover,
\begin{equation}\label{eq:Lambda-zero-expansion}
\Lambda(s)=-\frac1s+\frac{\gamma-\log(4\pi)}{2}+O(s),
\end{equation}
and
\begin{equation}\label{eq:xi-zero-expansion}
\xi(s)=\frac12+\frac{\log(4\pi)-2-\gamma}{4}s+O(s^2).
\end{equation}
In particular,
\[
\xi(0)=\xi(1)=\frac12.
\]
\end{proposition}

\begin{proof}
Start from \eqref{eq:functional-zeta}.  As $s\to 0$ we have
\[
2^s\pi^{s-1}=\pi^{-1}\bigl(1+\log(2\pi)s+O(s^2)\bigr),
\]
\[
\sin\!\left(\frac{\pi s}{2}\right)=\frac{\pi s}{2}+O(s^3),
\qquad
\Gamma(1-s)=1+\gamma s+O(s^2),
\]
and by Proposition~\ref{prop:stieltjes-expansion},
\[
\zeta(1-s)=-\frac1s+\gamma+O(s).
\]
Multiplying the last three expansions first gives
\[
\sin\!\left(\frac{\pi s}{2}\right)\Gamma(1-s)\zeta(1-s)
=\left(\frac{\pi s}{2}+O(s^3)\right)\left(1+\gamma s+O(s^2)\right)\left(-\frac1s+\gamma+O(s)\right)
=-\frac{\pi}{2}+O(s^2),
\]
because the linear terms cancel exactly.  Multiplying by $2^s\pi^{s-1}$ now yields
\[
\zeta(s)=\pi^{-1}\bigl(1+\log(2\pi)s+O(s^2)\bigr)\left(-\frac{\pi}{2}+O(s^2)\right)
=-\frac12-\frac12\log(2\pi)s+O(s^2),
\]
which is \eqref{eq:zeta-zero-expansion}.

Next, use the standard gamma expansion
\[
\Gamma\!\left(\frac{s}{2}\right)=\frac{2}{s}-\gamma+O(s)
\qquad (s\to 0)
\]
and
\[
\pi^{-s/2}=1-\frac12\log\pi\,s+O(s^2).
\]
Multiplying these with \eqref{eq:zeta-zero-expansion} gives
\[
\Lambda(s)=\left(1-\frac12\log\pi\,s+O(s^2)\right)
\left(\frac2s-\gamma+O(s)\right)
\left(-\frac12-\frac12\log(2\pi)s+O(s^2)\right)
=-\frac1s+\frac{\gamma-\log(4\pi)}{2}+O(s),
\]
which is \eqref{eq:Lambda-zero-expansion}.  Finally,
\[
\xi(s)=\frac12 s(s-1)\Lambda(s),
\]
so multiplying the previous expansion by $\frac12 s(s-1)$ gives \eqref{eq:xi-zero-expansion} and in particular $\xi(0)=1/2$.  The identity $\xi(1)=1/2$ follows from the functional equation $\xi(1-s)=\xi(s)$.
\end{proof}

\begin{proposition}\label{prop:zeta-negative-on-critical-strip}
For every real $s$ with $0<s<1$ one has
\[
\zeta(s)<0.
\]
Consequently,
\[
\xi\!\left(\frac12\right)>0.
\]
\end{proposition}

\begin{proof}
For $\Re(s)>0$, the Dirichlet eta function
\[
\eta(s):=\sum_{n\ge 1}\frac{(-1)^{n-1}}{n^s}
\]
satisfies
\[
\eta(s)=(1-2^{1-s})\zeta(s).
\]
If $0<s<1$ is real, then we may group the alternating series as
\[
\eta(s)=\sum_{k\ge 1}\left((2k-1)^{-s}-(2k)^{-s}\right),
\]
and every summand is strictly positive.  Hence $\eta(s)>0$.  On the same interval one has $1-2^{1-s}<0$, so dividing by this negative factor gives $\zeta(s)<0$.

At $s=1/2$ this shows $\zeta(1/2)<0$.  Since
\[
\xi\!\left(\frac12\right)=\frac12\cdot\frac12\cdot\left(-\frac12\right)\pi^{-1/4}\Gamma\!\left(\frac14\right)\zeta\!\left(\frac12\right)
=-\frac18\pi^{-1/4}\Gamma\!\left(\frac14\right)\zeta\!\left(\frac12\right),
\]
and the prefactor outside $\zeta(1/2)$ is positive, the second claim follows.
\end{proof}

\subsection{Trivial zeros and completed trivial weights}

\begin{proposition}\label{prop:trivial-derivative-and-Lambda}
For every integer $n\ge 1$ one has
\begin{equation}\label{eq:trivial-derivative-again}
\zeta'(-2n)=(-1)^n\frac{(2n)!}{2(2\pi)^{2n}}\,\zeta(2n+1),
\end{equation}
and therefore
\begin{equation}\label{eq:trivial-abs-derivative}
\bigl|\zeta'(-2n)\bigr|=\frac{(2n)!}{2(2\pi)^{2n}}\,\zeta(2n+1).
\end{equation}
In addition,
\begin{equation}\label{eq:Lambda-trivial-values}
\Lambda(-2n)=\Lambda(2n+1)=\frac{(2n)!}{4^n n!\,\pi^n}\,\zeta(2n+1),
\end{equation}
and the derivative and completed-value weights are related by
\begin{equation}\label{eq:trivial-weight-relation}
\bigl|\zeta'(-2n)\bigr|=\frac{n!}{2\pi^n}\,\Lambda(-2n).
\end{equation}
\end{proposition}

\begin{proof}
The derivative formula \eqref{eq:trivial-derivative-again} is obtained exactly as in the earlier split-zero note: write
\[
\zeta(s)=g_n(s)\sin\!\left(\frac{\pi s}{2}\right),
\qquad
g_n(s):=2^s\pi^{s-1}\Gamma(1-s)\zeta(1-s),
\]
and evaluate at $s=-2n$.  Since $g_n$ is holomorphic there and the sine factor has a simple zero, one gets
\[
\zeta'(-2n)=g_n(-2n)\cdot \frac{\pi}{2}\cos(-n\pi)
=(-1)^n\frac{(2n)!}{2(2\pi)^{2n}}\,\zeta(2n+1).
\]
Taking absolute values gives \eqref{eq:trivial-abs-derivative}.

For \eqref{eq:Lambda-trivial-values}, use the functional equation for the completed function:
\[
\Lambda(s)=\Lambda(1-s).
\]
Hence
\[
\Lambda(-2n)=\Lambda(2n+1)=\pi^{-(2n+1)/2}\Gamma\!\left(n+\frac12\right)\zeta(2n+1).
\]
The classical half-integer gamma identity
\[
\Gamma\!\left(n+\frac12\right)=\frac{(2n)!}{4^n n!}\sqrt\pi
\]
therefore gives
\[
\Lambda(-2n)=\frac{(2n)!}{4^n n!\,\pi^n}\,\zeta(2n+1).
\]
Finally, dividing \eqref{eq:trivial-abs-derivative} by \eqref{eq:Lambda-trivial-values} yields \eqref{eq:trivial-weight-relation}.
\end{proof}

\subsection{Leading jets at nontrivial zeros}

\begin{definition}
Let $\rho\in Z_{\mathrm{nt}}$ have multiplicity $m_\rho$.  We define the leading local jet coefficients
\[
\lambda_\rho^{\zeta}:=\frac{\zeta^{(m_\rho)}(\rho)}{m_\rho!},
\qquad
\lambda_\rho^{\xi}:=\frac{\xi^{(m_\rho)}(\rho)}{m_\rho!}.
\]
Equivalently,
\[
\zeta(s)=\lambda_\rho^{\zeta}(s-\rho)^{m_\rho}+O\bigl((s-\rho)^{m_\rho+1}\bigr),
\qquad
\xi(s)=\lambda_\rho^{\xi}(s-\rho)^{m_\rho}+O\bigl((s-\rho)^{m_\rho+1}\bigr).
\]
\end{definition}

\begin{proposition}\label{prop:leading-jets-nontrivial}
Let $\rho\in Z_{\mathrm{nt}}$.  Then
\begin{equation}\label{eq:jet-relation-zeta-xi}
\lambda_\rho^{\xi}=\frac12\rho(\rho-1)\pi^{-\rho/2}\Gamma\!\left(\frac{\rho}{2}\right)\lambda_\rho^{\zeta}.
\end{equation}
Moreover, if $m=m_\rho$, then the packet symmetries satisfy
\begin{equation}\label{eq:jet-symmetry-reflection}
\lambda_{1-\rho}^{\xi}=(-1)^m\lambda_\rho^{\xi},
\end{equation}
and
\begin{equation}\label{eq:jet-symmetry-conjugation}
\lambda_{\overline\rho}^{\xi}=\overline{\lambda_\rho^{\xi}}.
\end{equation}
Consequently, the absolute value
\begin{equation}\label{eq:packet-weight-def}
\omega(O):=\bigl|\lambda_\rho^{\xi}\bigr|
\end{equation}
is well defined on every packet $O\in Y_\xi$.
\end{proposition}

\begin{proof}
Write
\[
\phi(s):=\frac12 s(s-1)\pi^{-s/2}\Gamma\!\left(\frac{s}{2}\right).
\]
Since $\rho$ is a nontrivial zero, $\rho\neq 0,1,-2,-4,\dots$, so $\phi$ is holomorphic and nonzero at $\rho$.  Then $\xi=\phi\zeta$ near $\rho$.  If
\[
\zeta(s)=\lambda_\rho^{\zeta}(s-\rho)^m+O\bigl((s-\rho)^{m+1}\bigr),
\]
then multiplication by the holomorphic unit $\phi(s)=\phi(\rho)+O(s-\rho)$ gives
\[
\xi(s)=\phi(\rho)\lambda_\rho^{\zeta}(s-\rho)^m+O\bigl((s-\rho)^{m+1}\bigr),
\]
which is exactly \eqref{eq:jet-relation-zeta-xi}.

For the symmetry formulas, differentiate the identities
\[
\xi(1-s)=\xi(s),
\qquad
\overline{\xi(\overline s)}=\xi(s).
\]
The first gives
\[
(-1)^m\xi^{(m)}(1-s)=\xi^{(m)}(s),
\]
so evaluating at $s=\rho$ yields
\[
\xi^{(m)}(1-\rho)=(-1)^m\xi^{(m)}(\rho),
\]
which is \eqref{eq:jet-symmetry-reflection} after dividing by $m!$.  The second identity implies
\[
\overline{\xi^{(m)}(\overline s)}=\xi^{(m)}(s)
\]
for every $m$, hence \eqref{eq:jet-symmetry-conjugation}.  Formula \eqref{eq:packet-weight-def} is therefore independent of the choice of $\rho\in O$.
\end{proof}

\section{The centered entire function and the packet product}

\begin{definition}
Define the centered entire function
\[
\Xi(z):=\xi\!\left(\frac12+z\right).
\]
If $\rho$ is a nontrivial zero of $\zeta$, set
\[
\alpha_\rho:=\rho-\frac12.
\]
Then the nonzero zeros of $\Xi$ are exactly the points $\alpha_\rho$.
\end{definition}

\begin{lemma}\label{lem:Xi-even-and-real}
The function $\Xi$ is entire and even:
\[
\Xi(-z)=\Xi(z)
\qquad (z\in\CC).
\]
Moreover,
\[
\overline{\Xi(\overline z)}=\Xi(z),
\]
so the Taylor coefficients of $\Xi$ at $0$ are all real and every odd Taylor coefficient vanishes.
\end{lemma}

\begin{proof}
Entirety follows from the fact that $\xi$ is entire.  Evenness is immediate from the functional equation:
\[
\Xi(-z)=\xi\!\left(\frac12-z\right)=\xi\!\left(1-\left(\frac12+z\right)\right)=\xi\!\left(\frac12+z\right)=\Xi(z).
\]
Likewise,
\[
\overline{\Xi(\overline z)}=\overline{\xi\!\left(\frac12+\overline z\right)}=\xi\!\left(\frac12+z\right)=\Xi(z),
\]
using $\overline{\xi(\overline s)}=\xi(s)$.  Therefore the Taylor coefficients at $0$ are real, and evenness forces the odd coefficients to vanish.
\end{proof}

\begin{definition}
Let $\mathcal A$ be a multiset containing one representative from each sign pair $\{\alpha,-\alpha\}$ of nonzero zeros of $\Xi$, repeated according to multiplicity.  Thus the zero multiset of $\Xi$ is
\[
\{\pm\alpha:\alpha\in\mathcal A\},
\]
with multiplicities understood.
\end{definition}

\begin{theorem}[Centered sign-pair product]\label{thm:centered-sign-pair-product}
Assume the classical input that $\Xi$ is entire of order $1$.  Then the product
\begin{equation}\label{eq:Xi-sign-pair-product}
\Xi(z)=\Xi(0)\prod_{\alpha\in\mathcal A}\left(1-\frac{z^2}{\alpha^2}\right)
\end{equation}
converges absolutely and locally uniformly on compact subsets of $\CC$.  In particular,
\[
\sum_{\alpha\in\mathcal A}\frac{1}{|\alpha|^2}<\infty.
\]
\end{theorem}

\begin{proof}
By Hadamard factorization for order-one entire functions, one has
\[
\Xi(z)=e^{a+bz}\prod_{\beta\in Z(\Xi)\setminus\{0\}} E_1\!\left(\frac{z}{\beta}\right),
\qquad
E_1(w):=(1-w)e^w,
\]
with convergence in canonical product sense.  Because the nonzero zeros occur in sign pairs and $\Xi$ is even, we may group the factors pairwise:
\[
E_1\!\left(\frac{z}{\alpha}\right)E_1\!\left(-\frac{z}{\alpha}\right)
=(1-z/\alpha)e^{z/\alpha}(1+z/\alpha)e^{-z/\alpha}
=1-\frac{z^2}{\alpha^2}.
\]
Evenness also forces $b=0$: indeed,
\[
1=\frac{\Xi(z)}{\Xi(-z)}=e^{2bz}
\]
for all $z$, hence $b=0$.  Therefore
\[
\Xi(z)=e^a\prod_{\alpha\in\mathcal A}\left(1-\frac{z^2}{\alpha^2}\right).
\]
Evaluating at $z=0$ gives $e^a=\Xi(0)$, hence \eqref{eq:Xi-sign-pair-product}.

Finally, canonical product theory for an order-one entire function implies
\[
\sum_{\beta\in Z(\Xi)\setminus\{0\}} \frac{1}{|\beta|^2}<\infty,
\]
so after choosing one representative from each sign pair we still have
\[
\sum_{\alpha\in\mathcal A}\frac{1}{|\alpha|^2}<\infty.
\]
This summability implies that the ordinary product
\[
\prod_{\alpha\in\mathcal A}\left(1-\frac{z^2}{\alpha^2}\right)
\]
converges absolutely and locally uniformly on compact sets, because for $|z|\le R$ and $|\alpha|>2R$ one has
\[
\left|\frac{z^2}{\alpha^2}\right|\le \frac{R^2}{|\alpha|^2},
\]
and the tail is controlled by a summable majorant.
\end{proof}

\begin{definition}\label{def:packet-polynomial}
Let $Y_\xi=|D_{\mathrm{nt}}|/W$ be the packet quotient.  For a packet $O\in Y_\xi$, define its normalized packet polynomial by
\begin{equation}\label{eq:packet-polynomial-def}
P_O(s):=\prod_{\eta\in O}\frac{s-\eta}{\frac12-\eta}.
\end{equation}
This is well defined because $1/2\notin O$ by Proposition~\ref{prop:zeta-negative-on-critical-strip}.
\end{definition}

\begin{lemma}\label{lem:packet-polynomial-explicit}
For every packet $O\in Y_\xi$, the polynomial $P_O(s)$ has real coefficients and satisfies $P_O(1/2)=1$.

More precisely:
\begin{enumerate}[label=(\roman*)]
\item If $O=\{\frac12+i\gamma,\frac12-i\gamma\}$ is a critical-line packet, then
\begin{equation}\label{eq:packet-polynomial-line}
P_O(s)=1+\frac{(s-1/2)^2}{\gamma^2}.
\end{equation}
\item If $O=\{\beta\pm i\gamma,1-\beta\pm i\gamma\}$ with $\beta\neq 1/2$, then
\begin{equation}\label{eq:packet-polynomial-offline}
P_O(s)=\frac{\bigl((s-\beta)^2+\gamma^2\bigr)\bigl((s-1+\beta)^2+\gamma^2\bigr)}{\bigl((\beta-1/2)^2+\gamma^2\bigr)^2}.
\end{equation}
Equivalently, if $\alpha=(\beta-1/2)+i\gamma$, then
\begin{equation}\label{eq:packet-polynomial-offline-z}
P_O(s)=\left(1-\frac{(s-1/2)^2}{\alpha^2}\right)\left(1-\frac{(s-1/2)^2}{\overline\alpha^2}\right)
=1-2\Re(\alpha^{-2})(s-1/2)^2+|\alpha|^{-4}(s-1/2)^4.
\end{equation}
\end{enumerate}
\end{lemma}

\begin{proof}
The definition \eqref{eq:packet-polynomial-def} is manifestly normalized by $P_O(1/2)=1$.  Because each packet is stable under complex conjugation, the coefficients are real.

If $O=\{1/2\pm i\gamma\}$, then
\[
P_O(s)=\frac{(s-1/2-i\gamma)(s-1/2+i\gamma)}{(-i\gamma)(i\gamma)}
=\frac{(s-1/2)^2+\gamma^2}{\gamma^2}
=1+\frac{(s-1/2)^2}{\gamma^2}.
\]
This proves \eqref{eq:packet-polynomial-line}.

For an off-line packet,
\[
\prod_{\eta\in O}(s-\eta)=\bigl((s-\beta)^2+\gamma^2\bigr)\bigl((s-1+\beta)^2+\gamma^2\bigr).
\]
Likewise,
\[
\prod_{\eta\in O}\left(\frac12-\eta\right)=\bigl((\beta-1/2)^2+\gamma^2\bigr)^2.
\]
This gives \eqref{eq:packet-polynomial-offline}.  Writing $z=s-1/2$ and $\alpha=(\beta-1/2)+i\gamma$ transforms the same expression into \eqref{eq:packet-polynomial-offline-z}.
\end{proof}

\begin{theorem}[Packet product on the orbit space]\label{thm:packet-product-on-Y}
Let $m_O$ denote the common multiplicity of the points in the packet $O\in Y_\xi$.  Then
\begin{equation}\label{eq:packet-product-on-Y}
\xi(s)=\xi\!\left(\frac12\right)\prod_{O\in Y_\xi} P_O(s)^{m_O},
\end{equation}
where the product converges absolutely and locally uniformly on compact subsets of $\CC$.
\end{theorem}

\begin{proof}
Start from the centered sign-pair product \eqref{eq:Xi-sign-pair-product} with $z=s-1/2$:
\[
\xi(s)=\xi\!\left(\frac12\right)\prod_{\alpha\in\mathcal A}\left(1-\frac{(s-1/2)^2}{\alpha^2}\right).
\]
Conjugation preserves the multiset $\mathcal A$ up to permutation.  Group the factors according to conjugation orbits on $\mathcal A$.

If $\alpha=i\gamma$ is fixed by conjugation, the corresponding packet is critical-line and Lemma~\ref{lem:packet-polynomial-explicit}(i) shows that the factor is exactly $P_O(s)$.

If $\alpha$ and $\overline\alpha$ are distinct, they belong to the same packet $O$, and Lemma~\ref{lem:packet-polynomial-explicit}(ii) gives
\[
\left(1-\frac{(s-1/2)^2}{\alpha^2}\right)
\left(1-\frac{(s-1/2)^2}{\overline\alpha^2}\right)
=P_O(s).
\]
Because multiplicities are constant on packets, each packet contributes the power $P_O(s)^{m_O}$.  Thus \eqref{eq:packet-product-on-Y} holds.

Absolute and locally uniform convergence follow from Theorem~\ref{thm:centered-sign-pair-product}, because every packet factor is either one sign-pair factor or the product of a conjugate pair of sign-pair factors.
\end{proof}

\begin{corollary}[Packet formulation of RH]\label{cor:RH-packet-criterion}
The following are equivalent:
\begin{enumerate}[label=(\roman*)]
\item the Riemann hypothesis holds;
\item every packet $O\in Y_\xi$ has cardinality $2$;
\item in the product \eqref{eq:packet-product-on-Y}, every factor $P_O$ is quadratic of the form
\[
P_O(s)=1+\frac{(s-1/2)^2}{\gamma_O^2}
\qquad (\gamma_O>0),
\]
and no quartic packet factor occurs.
\end{enumerate}
\end{corollary}

\begin{proof}
A nontrivial zero lies on the critical line if and only if its packet under $W$ has two points rather than four.  Thus (i) and (ii) are equivalent.  By Lemma~\ref{lem:packet-polynomial-explicit}, cardinality-$2$ packets produce precisely the quadratic factors, while cardinality-$4$ packets produce the real quartic factors.  Hence (ii) and (iii) are equivalent.
\end{proof}

\section{Centered Taylor coefficients and symmetric-function formulas}

\begin{definition}
Choose an ordering $\mathcal A=\{\alpha_1,\alpha_2,\dots\}$ of the multiset from Theorem~\ref{thm:centered-sign-pair-product}, and set
\[
a_j:=\alpha_j^{-2}.
\]
Because
\[
\sum_{j\ge 1}|a_j|<\infty,
\]
all elementary symmetric sums
\begin{equation}\label{eq:elementary-symmetric-def}
\sigma_n:=\sum_{1\le j_1<\cdots<j_n} a_{j_1}\cdots a_{j_n}
\qquad (n\ge 1),
\end{equation}
are absolutely convergent.  We also set $\sigma_0:=1$.
\end{definition}

\begin{lemma}\label{lem:infinite-symmetric-expansion}
For every $z\in\CC$ one has the absolutely convergent expansion
\begin{equation}\label{eq:infinite-symmetric-expansion}
\prod_{j\ge 1}(1-a_j z^2)=\sum_{n=0}^{\infty}(-1)^n\sigma_n z^{2n},
\end{equation}
and the convergence is locally uniform on compact subsets of $\CC$.
\end{lemma}

\begin{proof}
Fix $R>0$ and consider $|z|\le R$.  Since $\sum_j |a_j|<\infty$, one has
\[
\sum_j |a_j|R^2<\infty.
\]
Therefore the infinite product
\[
\prod_j (1-a_j z^2)
\]
converges absolutely and uniformly on $|z|\le R$.  For each finite $N$,
\[
\prod_{j=1}^N(1-a_j z^2)=\sum_{n=0}^{N}(-1)^n\sigma_n^{(N)} z^{2n},
\]
where $\sigma_n^{(N)}$ is the $n$th elementary symmetric sum in $a_1,\dots,a_N$.  For fixed $n$, the coefficients $\sigma_n^{(N)}$ converge absolutely to $\sigma_n$ as $N\to\infty$, because
\[
\sum_{1\le j_1<\cdots<j_n<\infty}|a_{j_1}\cdots a_{j_n}|
\le \frac{1}{n!}\left(\sum_{j\ge 1}|a_j|\right)^n<\infty.
\]
Passing to the limit in the finite expansions yields \eqref{eq:infinite-symmetric-expansion}.  Uniform convergence on $|z|\le R$ follows from the Weierstrass $M$-test, since
\[
\sum_{n\ge 0}|\sigma_n|R^{2n}\le \prod_{j\ge 1}(1+|a_j|R^2)<\infty.
\]
\end{proof}

\begin{theorem}[Exact Taylor coefficients of the centered completed zeta]\label{thm:Xi-taylor-coefficients}
The centered function $\Xi$ has the exact Taylor expansion
\begin{equation}\label{eq:Xi-taylor-symmetric}
\Xi(z)=\Xi(0)\sum_{n=0}^{\infty}(-1)^n\sigma_n z^{2n},
\end{equation}
where the coefficients $\sigma_n$ are the absolutely convergent elementary symmetric sums from \eqref{eq:elementary-symmetric-def}.  Equivalently,
\begin{equation}\label{eq:Xi-even-derivatives}
\frac{\Xi^{(2n)}(0)}{(2n)!\,\Xi(0)}=(-1)^n\sigma_n,
\qquad
\Xi^{(2n+1)}(0)=0
\qquad (n\ge 0).
\end{equation}
Translating back to $s$ gives
\begin{equation}\label{eq:xi-central-derivatives}
\frac{\xi^{(2n)}(1/2)}{(2n)!\,\xi(1/2)}=(-1)^n\sigma_n,
\qquad
\xi^{(2n+1)}(1/2)=0.
\end{equation}
In particular,
\begin{align}
\frac{\xi''(1/2)}{2\,\xi(1/2)}&=-\sum_{j\ge 1}\frac{1}{\alpha_j^2},\label{eq:xi-second-derivative}\\
\frac{\xi^{(4)}(1/2)}{24\,\xi(1/2)}&=\sum_{1\le i<j<\infty}\frac{1}{\alpha_i^2\alpha_j^2}.\label{eq:xi-fourth-derivative}
\end{align}
\end{theorem}

\begin{proof}
Insert the product formula \eqref{eq:Xi-sign-pair-product} into Lemma~\ref{lem:infinite-symmetric-expansion}.  This gives \eqref{eq:Xi-taylor-symmetric}.  Comparing coefficients with the Taylor series of $\Xi$ yields \eqref{eq:Xi-even-derivatives}.  Since $\Xi(z)=\xi(1/2+z)$, the identities \eqref{eq:xi-central-derivatives} follow immediately.  The last two displayed formulas are the cases $n=1$ and $n=2$.
\end{proof}

\begin{remark}
The coefficients in Theorem~\ref{thm:Xi-taylor-coefficients} are the exact analytic counterpart of the symmetric-function expansions that appeared elsewhere in the uploaded audit note.  Here the variables are no longer cyclotomic affine factors but the reciprocal squares of the centralized packet locations.
\end{remark}

\section{Global pole-trivial-packet factorization of $\zeta$}

\subsection{Removing the trivial zeros with the reciprocal gamma product}

\begin{theorem}[Weierstrass product for the trivial sector]\label{thm:gamma-product}
For every $s\in\CC$ one has the locally uniformly convergent product
\begin{equation}\label{eq:gamma-weierstrass}
\frac{1}{\Gamma(s/2)}=\frac{s}{2}e^{\gamma s/2}\prod_{n\ge 1}\left(1+\frac{s}{2n}\right)e^{-s/(2n)}.
\end{equation}
Consequently,
\begin{equation}\label{eq:Lambda-to-zeta-via-gamma-product}
\zeta(s)=\frac{\pi^{s/2}e^{\gamma s/2}}{s-1}\left(\prod_{n\ge 1}\left(1+\frac{s}{2n}\right)e^{-s/(2n)}\right)\xi(s).
\end{equation}
\end{theorem}

\begin{proof}
Equation \eqref{eq:gamma-weierstrass} is the classical Weierstrass product for the reciprocal gamma function, applied with $z=s/2$.

For the second formula, solve the defining relation
\[
\xi(s)=\frac12 s(s-1)\pi^{-s/2}\Gamma\!\left(\frac{s}{2}\right)\zeta(s)
\]
for $\zeta(s)$:
\[
\zeta(s)=\frac{2\pi^{s/2}}{s(s-1)}\,\frac{\xi(s)}{\Gamma(s/2)}.
\]
Substituting \eqref{eq:gamma-weierstrass} into this identity cancels the factor $2/s$ and yields \eqref{eq:Lambda-to-zeta-via-gamma-product}.
\end{proof}

\subsection{The full packet factorization}

\begin{theorem}[Pole-trivial-packet factorization of the Riemann zeta function]\label{thm:full-zeta-factorization}
The Riemann zeta function admits the exact factorization
\begin{equation}\label{eq:full-zeta-factorization}
\zeta(s)
=
\frac{\xi(1/2)e^{(\gamma+\log\pi)s/2}}{s-1}
\prod_{n\ge 1}\left(1+\frac{s}{2n}\right)e^{-s/(2n)}
\prod_{O\in Y_\xi} P_O(s)^{m_O},
\end{equation}
where $P_O$ is the normalized packet polynomial from Definition~\ref{def:packet-polynomial} and $m_O$ is the common multiplicity on the packet $O$.  The product over $Y_\xi$ converges absolutely and locally uniformly on compact subsets of $\CC$.

Thus $\zeta$ decomposes canonically into three factors:
\begin{enumerate}[label=(\roman*)]
\item the pole factor $1/(s-1)$;
\item the trivial-zero factor
\[
e^{(\gamma+\log\pi)s/2}\prod_{n\ge 1}\left(1+\frac{s}{2n}\right)e^{-s/(2n)};
\]
\item the nontrivial packet factor
\[
\prod_{O\in Y_\xi} P_O(s)^{m_O}.
\]
\end{enumerate}
\end{theorem}

\begin{proof}
Insert the packet product \eqref{eq:packet-product-on-Y} into \eqref{eq:Lambda-to-zeta-via-gamma-product}.  Since
\[
\pi^{s/2}e^{\gamma s/2}=e^{(\gamma+\log\pi)s/2},
\]
this gives exactly \eqref{eq:full-zeta-factorization}.  The convergence statement follows from Theorem~\ref{thm:packet-product-on-Y}.
\end{proof}

\begin{remark}
The factorization \eqref{eq:full-zeta-factorization} is the precise sense in which the entire nontrivial-zero sector is ``carried'' by the packet quotient, while the trivial-zero sector is carried by the reciprocal gamma factor.  In the split-zero divisor language, the earlier Boolean shadow now receives a full analytic factor over each packet.
\end{remark}

\section{Packet-grouped logarithmic derivatives}

\begin{proposition}\label{prop:packet-log-derivative}
For every packet $O\in Y_\xi$ one has
\begin{equation}\label{eq:packet-log-derivative}
\frac{P_O'(s)}{P_O(s)}=\sum_{\eta\in O}\frac{1}{s-\eta}.
\end{equation}
In particular, the right-hand side is a rational function with real coefficients.
\end{proposition}

\begin{proof}
By definition,
\[
P_O(s)=\prod_{\eta\in O}\frac{s-\eta}{1/2-\eta}.
\]
Taking the logarithmic derivative annihilates the constant denominator and yields
\[
\frac{P_O'(s)}{P_O(s)}=\sum_{\eta\in O}\frac{1}{s-\eta}.
\]
Because $P_O$ has real coefficients by Lemma~\ref{lem:packet-polynomial-explicit}, the rational function on the left does as well.
\end{proof}

\begin{theorem}[Packet-grouped partial-fraction expansion of $\zeta'/\zeta$]\label{thm:zeta-log-derivative-packet}
On the domain
\[
\CC\setminus\Bigl(\{1\}\cup\{-2,-4,-6,\dots\}\cup |D_{\mathrm{nt}}|\Bigr)
\]
one has the exact identity
\begin{equation}\label{eq:zeta-log-derivative-packet}
\frac{\zeta'(s)}{\zeta(s)}
=
\frac{\gamma+\log\pi}{2}
-\frac{1}{s-1}
+\sum_{n\ge 1}\left(\frac{1}{s+2n}-\frac{1}{2n}\right)
+\sum_{O\in Y_\xi} m_O\frac{P_O'(s)}{P_O(s)}.
\end{equation}
Equivalently,
\begin{equation}\label{eq:zeta-log-derivative-packet-expanded}
\frac{\zeta'(s)}{\zeta(s)}
=
\frac{\gamma+\log\pi}{2}
-\frac{1}{s-1}
+\sum_{n\ge 1}\left(\frac{1}{s+2n}-\frac{1}{2n}\right)
+\sum_{O\in Y_\xi} m_O\sum_{\eta\in O}\frac{1}{s-\eta}.
\end{equation}
Every pole on the right-hand side comes from exactly one of the three sectors: the pole at $1$, the trivial zeros, or the nontrivial packets.
\end{theorem}

\begin{proof}
Differentiate logarithmically the factorization \eqref{eq:full-zeta-factorization}.  The constant factor $\xi(1/2)$ disappears, while
\[
\frac{d}{ds}\left(\frac{\gamma+\log\pi}{2}s\right)=\frac{\gamma+\log\pi}{2},
\qquad
\frac{d}{ds}\log\frac{1}{s-1}=-\frac{1}{s-1},
\]
and
\[
\frac{d}{ds}\log\left(\left(1+\frac{s}{2n}\right)e^{-s/(2n)}\right)
=\frac{1}{s+2n}-\frac{1}{2n}.
\]
Finally, the packet factor contributes
\[
\sum_{O\in Y_\xi} m_O\frac{P_O'(s)}{P_O(s)}.
\]
This proves \eqref{eq:zeta-log-derivative-packet}.  Formula \eqref{eq:zeta-log-derivative-packet-expanded} is Proposition~\ref{prop:packet-log-derivative}.
\end{proof}

\begin{corollary}\label{cor:xi-log-derivative-packet}
On $\CC\setminus |D_{\mathrm{nt}}|$ one has
\begin{equation}\label{eq:xi-log-derivative-packet}
\frac{\xi'(s)}{\xi(s)}=\sum_{O\in Y_\xi} m_O\frac{P_O'(s)}{P_O(s)}
=\sum_{O\in Y_\xi} m_O\sum_{\eta\in O}\frac{1}{s-\eta}.
\end{equation}
\end{corollary}

\begin{proof}
Differentiate logarithmically the packet factorization \eqref{eq:packet-product-on-Y}.  The constant $\xi(1/2)$ disappears, and Proposition~\ref{prop:packet-log-derivative} gives the second equality.
\end{proof}

\section{Analytic decoration of the split-zero packet sheaf}

The earlier split-zero note descended the Boolean shadow of the nontrivial-zero divisor to the packet sheaf
\[
\mathcal P_\xi^{\cV}(U)=\prod_{O\in U} A^{m_O},
\qquad U\subset Y_\xi.
\]
The results above show that this packet sheaf carries two additional canonical analytic decorations.

\begin{definition}
For each packet $O\in Y_\xi$, choose any $\rho\in O$ and define
\[
\omega(O):=\bigl|\lambda_\rho^{\xi}\bigr|.
\]
By Proposition~\ref{prop:leading-jets-nontrivial}, this is well defined.
\end{definition}

\begin{definition}
For each packet $O\in Y_\xi$, let $P_O(s)$ be the normalized packet polynomial from Definition~\ref{def:packet-polynomial}.  We call the triple
\begin{equation}\label{eq:analytic-packet-data}
\bigl(m_O,\,P_O,\,\omega(O)\bigr)
\end{equation}
the \emph{full analytic packet datum} of $O$.
\end{definition}

\begin{proposition}\label{prop:analytic-packet-data-complete}
The packet datum \eqref{eq:analytic-packet-data} controls simultaneously the local, descended, and global contributions of the packet $O$:
\begin{enumerate}[label=(\roman*)]
\item $m_O$ is the multiplicity of the packet in the Boolean shadow stalk $A^{m_O}$;
\item $\omega(O)$ is the absolute size of the leading local jet of $\xi$ on the packet;
\item $P_O(s)^{m_O}$ is the exact normalized global holomorphic factor contributed by $O$ to \eqref{eq:packet-product-on-Y} and hence to \eqref{eq:full-zeta-factorization};
\item the rational function $m_O P_O'(s)/P_O(s)$ is the exact contribution of $O$ to the logarithmic derivative \eqref{eq:zeta-log-derivative-packet}.
\end{enumerate}
\end{proposition}

\begin{proof}
Part (i) is built into the definition of the descended Boolean packet sheaf.  Part (ii) is the packet invariance of Proposition~\ref{prop:leading-jets-nontrivial}.  Part (iii) is exactly the packet product theorem, and part (iv) is Proposition~\ref{prop:packet-log-derivative} applied inside Theorem~\ref{thm:zeta-log-derivative-packet}.
\end{proof}
% END part_analytic

% BEGIN part_connes
\section{The spectral plane and the Connes--Consani realization}

\subsection{The affine spectral coordinate}

\begin{definition}
Define the affine map
\[
\Phi:\CC\longrightarrow\CC,
\qquad
\Phi(s):=i\Bigl(s-\frac12\Bigr).
\]
Its inverse is
\[
\Phi^{-1}(z)=\frac12-iz.
\]
\end{definition}

\begin{proposition}\label{prop:Phi-symmetry}
For all $s\in\CC$ one has
\[
\Phi(1-s)=-\Phi(s),
\qquad
\Phi(\bar s)=-\overline{\Phi(s)}.
\]
Consequently $\Phi$ conjugates the action of $W$ on the $s$-plane to the action of the Klein four-group
\[
\widetilde W:=\{\mathrm{id},\sigma,\kappa,\sigma\kappa\}
\]
on the $z$-plane, where
\[
\sigma(z):=-z,
\qquad
\kappa(z):=-\bar z.
\]
In particular, if $\rho=\beta+i\gamma$, then
\[
\Phi(\rho)=-\gamma+i\Bigl(\beta-\frac12\Bigr),
\]
and the packet
\[
\{\rho,\bar\rho,1-\rho,1-\bar\rho\}
\]
is sent to
\[
\{z,-z,\bar z,-\bar z\},
\qquad z=\Phi(\rho).
\]
\end{proposition}

\begin{proof}
The identities are direct computations:
\[
\Phi(1-s)=i\Bigl(1-s-\frac12\Bigr)=i\Bigl(\frac12-s\Bigr)=-i\Bigl(s-\frac12\Bigr)=-\Phi(s),
\]
and
\[
\Phi(\bar s)=i\Bigl(\bar s-\frac12\Bigr)=-\overline{i\Bigl(s-\frac12\Bigr)}=-\overline{\Phi(s)}.
\]
Thus $w$ is conjugated to $\sigma$ and $c$ is conjugated to $\kappa$.  Applying these formulas to the four points of the $W$-orbit of $\rho$ yields the displayed spectral packet.
\end{proof}

\begin{definition}
Let
\[
Z_E:=\Phi(Z_\xi),
\qquad
D_E:=\Phi_*D_\xi=\sum_{z\in Z_E} m_z[z],
\\\qquad m_z:=m_{\Phi^{-1}(z)}.
\]
Let
\[
Y_E:=Z_E/\widetilde W
\]
be the spectral packet space.
\end{definition}

\begin{corollary}\label{cor:packet-space-iso}
The map $\Phi$ induces an isomorphism of discrete spaces
\[
\bar\Phi:Y_\xi\stackrel{\sim}{\longrightarrow} Y_E.
\]
\end{corollary}

\begin{proof}
By Proposition \ref{prop:Phi-symmetry}, $\Phi$ is equivariant for the isomorphism $W\cong\widetilde W$.  Since $\Phi$ is bijective on $\CC$, it restricts to a bijection $Z_\xi\to Z_E$ and therefore induces a bijection of quotient spaces.  Both spaces are discrete, so the induced bijection is a homeomorphism.
\end{proof}

\subsection{The Connes--Consani spectral divisor}

We now connect $D_E$ to the spectral realization of Connes and Consani.  Let $\cS(\RR)^{\mathrm{ev}}$ be the even Schwartz space on $\RR$, and let
\[
\cS(\RR)^{\mathrm{ev}}_1:=\left\{f\in\cS(\RR)^{\mathrm{ev}}: \int_{\RR} f(x)\,dx=0\right\}.
\]
For $f\in\cS(\RR)^{\mathrm{ev}}$, set
\[
(Ef)(u):=u^{1/2}\sum_{n\ge1} f(nu)
\qquad (u\in\RR_{>0}),
\]
and define its Mellin--Fourier transform by
\[
F(Ef)(z):=\int_{\RR_{>0}} (Ef)(u)u^{-iz}\,d^*u,
\qquad d^*u:=\frac{du}{u}.
\]

\begin{theorem}[Connes--Consani, \cite{ConnesConsaniCycles}]\label{thm:Connes-CC}
For every $f\in \cS(\RR)^{\mathrm{ev}}_1$ one has
\[
F(Ef)(z)=\zeta\Bigl(\frac12-iz\Bigr)\psi_f(z)
\]
on the half-plane $\Ima(z)>-\frac12$, where $\psi_f$ is holomorphic.  Consequently the common zero set of the functions $F(Ef)$, $f\in\cS(\RR)^{\mathrm{ev}}_1$, in the strip $|\Ima(z)|<\frac12$ is exactly the zero set of
\[
z\longmapsto \zeta\Bigl(\frac12-iz\Bigr).
\]
Moreover, the operator $\Delta=H(1+H)$ of \cite[Section~4]{ConnesConsaniCycles} has spectrum
\[
\left\{-\rho(1-\rho):\rho\in\CC\setminus\RR,\ \zeta(\rho)=0\right\}
\]
counted with multiplicity on the quotient described in \cite[Proposition~4.2]{ConnesConsaniCycles}.
\end{theorem}

\begin{proof}
This is exactly the content of equation~(2.2) and the common-zero statement in Section~2, together with Proposition~4.2 and its proof in Section~4 of \cite{ConnesConsaniCycles}.  The spectral values are written there as
\[
\Bigl(\rho-\frac12\Bigr)^2-\frac14,
\]
which is equal to $-\rho(1-\rho)$.
\end{proof}

\begin{theorem}\label{thm:bridge}
The affine map $\Phi$ identifies $Z_\xi$ with the Connes--Consani spectral support:
\[
Z_E=\left\{z\in\CC:\zeta\Bigl(\frac12-iz\Bigr)=0\right\}.
\]
Equivalently, $Z_E$ is the common zero set of the family $F(Ef)$, $f\in\cS(\RR)^{\mathrm{ev}}_1$, in the strip $|\Ima(z)|<\frac12$.

For every $\rho\in Z_\xi$ one has
\[
-\rho(1-\rho)=-\Phi(\rho)^2-\frac14.
\]
Thus $\Phi(\rho)$ is the square-root spectral coordinate for the Connes Laplacian.
\end{theorem}

\begin{proof}
Let $\rho\in Z_\xi$, so by definition of $D_\xi$ it is a non-trivial zero of $\zeta$.  Put $z=\Phi(\rho)$.  Since $\Phi^{-1}(z)=\frac12-iz$, one has
\[
\zeta(\rho)=0
\iff
\zeta\Bigl(\frac12-iz\Bigr)=0.
\]
Hence
\[
Z_E=\left\{z\in\CC:\zeta\Bigl(\frac12-iz\Bigr)=0\right\}.
\]
The strip condition is automatic, because if $\rho=\beta+i\gamma$ is non-trivial then $0<\beta<1$, and therefore
\[
\Ima(\Phi(\rho))=\beta-\frac12
\]
satisfies $|\Ima(\Phi(\rho))|<\frac12$.

The common-zero statement now follows from Theorem \ref{thm:Connes-CC}.

For the Laplacian relation, a direct calculation gives
\[
-\Phi(\rho)^2-\frac14
=-\Bigl(i\Bigl(\rho-\frac12\Bigr)\Bigr)^2-\frac14
=\Bigl(\rho-\frac12\Bigr)^2-\frac14
=-\rho(1-\rho).
\]
This is exactly the spectral parameter of $\Delta$ attached to $\rho$ by Theorem \ref{thm:Connes-CC}.
\end{proof}

\begin{definition}
The \emph{spectral Boolean packet sheaf} on $Y_E$ is
\[
\cP_E(U):=\prod_{P\in U} A^{m_P}
\qquad (U\subset Y_E\ \text{open}),
\]
where $m_P$ is the common multiplicity of the elements of the packet $P$.
\end{definition}

\begin{corollary}\label{cor:packet-transport}
Via the isomorphism $\bar\Phi:Y_\xi\to Y_E$, the packet sheaf $\cP_\xi$ transports canonically to $\cP_E$:
\[
\bar\Phi_*\cP_\xi\cong \cP_E.
\]
\end{corollary}

\begin{proof}
By definition of $D_E$, multiplicities are preserved by $\Phi$.  Hence for every open set $U\subset Y_E$ one has
\[
(\bar\Phi_*\cP_\xi)(U)=\cP_\xi(\bar\Phi^{-1}(U))
=\prod_{O\in \bar\Phi^{-1}(U)} A^{m_O}
=\prod_{P\in U} A^{m_P}
=\cP_E(U).
\]
These identifications commute with restriction maps.
\end{proof}

\section{Critical packets and $\zeta$-cycles}

\begin{proposition}\label{prop:critical-real}
Let $\rho\in Z_\xi$ and let $z=\Phi(\rho)$.  Then
\[
\Rea(\rho)=\frac12
\iff
z\in\RR.
\]
If $\rho=\frac12+i\gamma$, then
\[
\Phi(\rho)=-\gamma.
\]
Consequently the spectral packet of a critical zero is the two-point set
\[
\{-\gamma,\gamma\}.
\]
\end{proposition}

\begin{proof}
Write $\rho=\beta+i\gamma$.  Then Proposition \ref{prop:Phi-symmetry} gives
\[
\Phi(\rho)=-\gamma+i\Bigl(\beta-\frac12\Bigr).
\]
Thus $\Phi(\rho)$ is real if and only if $\beta=\frac12$.  In that case $\Phi(\rho)=-\gamma$.  The packet formula of Proposition \ref{prop:Phi-symmetry} then yields the two-point set $\{-\gamma,\gamma\}$.
\end{proof}

\begin{definition}
Let $Y_E^{\mathrm{crit}}\subset Y_E$ be the subset of critical packets, i.e. the packets contained in $\RR$.  For $P\in Y_E^{\mathrm{crit}}$, define
\[
\nu(P):=|z|
\]
for any $z\in P$.  This is well-defined because $P=\{z,-z\}$.

Define also the \emph{primitive cycle length}
\[
\Lambda(P):=\frac{2\pi}{\nu(P)}.
\]
\end{definition}

\begin{theorem}\label{thm:zeta-cycles}
Let $P\in Y_E^{\mathrm{crit}}$ and write $\nu(P)=s>0$.  Then every circle whose length is an integral multiple of
\[
\frac{2\pi}{s}
\]
is a $\zeta$-cycle in the sense of Connes and Consani.
\end{theorem}

\begin{proof}
Let $P=\{\pm s\}$ with $s>0$.  By Proposition \ref{prop:critical-real}, $P$ comes from a critical zero
\[
\rho=\frac12+is
\]
of $\zeta$.  Theorem~1.1(ii) of Connes and Consani, quoted in \cite{ConnesConsaniCycles}, states that if
\[
\zeta\Bigl(\frac12+is\Bigr)=0,
\]
then any circle whose length is an integral multiple of $2\pi/s$ is a $\zeta$-cycle.  Applying this with the displayed critical zero proves the assertion.
\end{proof}

\begin{corollary}\label{cor:critical-map}
The assignment
\[
\Lambda:Y_E^{\mathrm{crit}}\longrightarrow \RR_{>0},
\qquad
P\longmapsto \frac{2\pi}{\nu(P)},
\]
is a well-defined map from critical packets to primitive $\zeta$-cycle lengths.
\end{corollary}

\begin{proof}
This is immediate from the definition of $\nu(P)$ and Theorem \ref{thm:zeta-cycles}.
\end{proof}

\begin{remark}
The critical part of the packet map is therefore geometric in the strict Connes--Consani sense: besides the Boolean coefficient object, a critical zero packet determines a primitive circle length for the $\zeta$-cycles on the Scaling Site, together with all of its covering multiples.
\end{remark}
% END part_connes

% BEGIN part_phase
\section{Atomic phase factors and orthogonal walks}

The exact phase formalism used below is extracted from Tobley's unpublished correspondence on arctangent decompositions, orthogonal walks, and atomic step models of complex multiplication \cite{Tobley}.

\begin{definition}
For $a\in\RR$, define the normalized affine phase factor
\[
u(a):=\frac{1+\ii a}{\sqrt{1+a^2}}\in\Sph.
\]
For a finite sequence $a_1,\dots,a_M\in\RR$, define the unnormalized and normalized products
\[
U_M(a_1,\dots,a_M):=\prod_{n=1}^M (1+\ii a_n),
\qquad
\widehat U_M(a_1,\dots,a_M):=\prod_{n=1}^M \nu(a_n).
\]
\end{definition}

\begin{proposition}[Finite orthogonal-walk expansion]\label{prop:finite-orthogonal}
Let $a_1,\dots,a_M\in\RR$, and let $e_r(a_1,\dots,a_M)$ be the elementary symmetric polynomials.  Then
\[
U_M(a_1,\dots,a_M)=\sum_{r=0}^M \ii^r e_r(a_1,\dots,a_M).
\]
Consequently, if
\[
S_r:=\sum_{q=0}^r \ii^q e_q(a_1,\dots,a_M),
\qquad 0\le r\le M,
\]
then the successive differences $S_r-S_{r-1}$ are alternately horizontal and vertical segments, and the $r$th segment has signed length $e_r(a_1,\dots,a_M)$.

Moreover,
\[
\widehat U_M(a_1,\dots,a_M)=\exp\Bigl(\ii\sum_{n=1}^M \arctan a_n\Bigr).
\]
\end{proposition}

\begin{proof}
Expanding the product gives
\[
\prod_{n=1}^M (1+\ii a_n)=\sum_{I\subseteq [M]} \ii^{|I|}\prod_{j\in I} a_j.
\]
Grouping by the cardinality $|I|=r$ yields the first formula because
\[
e_r(a_1,\dots,a_M)=\sum_{|I|=r}\prod_{j\in I}a_j.
\]
Since multiplication by $\ii^r$ rotates the real axis through successive quarter-turns, the geometric statement is immediate.

For the normalized identity, note that
\[
\nu(a)=\frac{1+\ii a}{\sqrt{1+a^2}}=
\cos(\arctan a)+\ii\sin(\arctan a)=\e^{\ii\arctan a}.
\]
Multiplying over $n$ proves the claim.
\end{proof}

\begin{proposition}[Infinitesimal orthogonal rotation]\label{prop:derivative-nu}
For $a\in\RR$ one has
\[
\nu'(a)=\frac{\ii}{1+a^2}\,\nu(a).
\]
Consequently, if $a_n:I\to\RR$ are differentiable on an interval $I$, then
\[
\frac{d}{dt}\widehat U_M\bigl(a_1(t),\dots,a_M(t)\bigr)
=\ii\Bigl(\sum_{n=1}^M \frac{a_n'(t)}{1+a_n(t)^2}\Bigr)
\widehat U_M\bigl(a_1(t),\dots,a_M(t)\bigr).
\]
\end{proposition}

\begin{proof}
By Proposition~\ref{prop:finite-orthogonal},
\[
\nu(a)=\e^{\ii\arctan a},
\]
so
\[
\nu'(a)=\ii\frac{1}{1+a^2}\e^{\ii\arctan a}=\frac{\ii}{1+a^2}\,\nu(a).
\]
For the product, differentiate termwise and factor out the common product.
\end{proof}

\begin{proposition}[Equal-step limit]\label{prop:equal-step-limit}
For every $\theta\in\RR$,
\[
\lim_{N\to\infty}\nu\!\left(\frac{\theta}{N}\right)^N=\e^{\ii\theta}.
\]
\end{proposition}

\begin{proof}
By Proposition~\ref{prop:finite-orthogonal},
\[
\nu\!\left(\frac{\theta}{N}\right)^N
=\exp\!\Bigl(\ii N\arctan\!\left(\frac{\theta}{N}\right)\Bigr).
\]
Since $N\arctan(\theta/N)\to \theta$, the conclusion follows.
\end{proof}

\begin{theorem}[Exact Taylor atomization]\label{thm:taylor-atomization}
Let
\[
\Theta(t)=\sum_{n=1}^{\infty}\beta_n t^n
\]
be a real-analytic function on $(-R,R)$ with $\Theta(0)=0$ and $\beta_n\in\RR$.  Then there exists $r\in(0,R)$ such that for every $|t|\le r$ and every $n\ge1$ one has
\[
|\beta_n t^n|<\frac{\pi}{4}.
\]
Define the atomic slopes
\[
\alpha_n(t):=\tan\bigl(\beta_n t^n\bigr),\qquad |t|\le r.
\]
Then the following hold.
\begin{enumerate}[label=(\roman*)]
\item The series $\sum_{n\ge1}|\alpha_n(t)|$ converges uniformly on $[-r,r]$.
\item For every $M\ge1$,
\[
\prod_{n=1}^M \nu\bigl(\alpha_n(t)\bigr)=\exp\!\Bigl(\ii\sum_{n=1}^M\beta_n t^n\Bigr).
\]
\item The infinite product converges uniformly on $[-r,r]$, and
\[
\prod_{n=1}^{\infty}\nu\bigl(\alpha_n(t)\bigr)=\e^{\ii\Theta(t)}.
\]
\item The unnormalized finite products admit the orthogonal-walk expansions
\[
\prod_{n=1}^M (1+\ii\alpha_n(t))=
\sum_{q=0}^M \ii^q e_q\bigl(\alpha_1(t),\dots,\alpha_M(t)\bigr).
\]
Hence the normalized partial products are normalized endpoints of finite orthogonal walks with step lengths $e_q\bigl(\alpha_1(t),\dots,\alpha_M(t)\bigr)$.
\end{enumerate}
\end{theorem}

\begin{proof}
Choose $r\in(0,R)$ so small that
\[
\sum_{n=1}^{\infty}|\beta_n|r^n<\frac{\pi}{4}.
\]
Then for every $n$ and every $|t|\le r$,
\[
|\beta_n t^n|\le \sum_{m=1}^{\infty}|\beta_m|r^m<\frac{\pi}{4}.
\]
On $[-\pi/4,\pi/4]$ one has $|\tan x|\le 2|x|$, so
\[
|\alpha_n(t)|=|\tan(\beta_n t^n)|\le 2|\beta_n|r^n.
\]
Since $\sum_n |\beta_n|r^n<\infty$, the series $\sum_n |\alpha_n(t)|$ converges uniformly on $[-r,r]$.  This proves~(i).

For~(ii), observe that $|\beta_n t^n|<\pi/2$, so
\[
\arctan\alpha_n(t)=\arctan\bigl(\tan(\beta_n t^n)\bigr)=\beta_n t^n.
\]
By Proposition~\ref{prop:finite-orthogonal},
\[
\prod_{n=1}^M \nu\bigl(\alpha_n(t)\bigr)
=\exp\!\Bigl(\ii\sum_{n=1}^M \arctan\alpha_n(t)\Bigr)
=\exp\!\Bigl(\ii\sum_{n=1}^M \beta_n t^n\Bigr).
\]
This proves~(ii).

Since the power series for $\Theta$ converges uniformly on $[-r,r]$, the partial sums
\[
\Theta_M(t):=\sum_{n=1}^M\beta_n t^n
\]
converge uniformly to $\Theta(t)$.  By~(ii),
\[
\prod_{n=1}^M \nu\bigl(\alpha_n(t)\bigr)=\e^{\ii\Theta_M(t)},
\]
so the partial products converge uniformly to $\e^{\ii\Theta(t)}$.  This proves~(iii).

Part~(iv) is Proposition~\ref{prop:finite-orthogonal} applied to the finite sequence $\alpha_1(t),\dots,\alpha_M(t)$.
\end{proof}

\begin{corollary}[Two exact walk models for the same phase]\label{cor:two-walks}
Under the hypotheses of Theorem~\ref{thm:taylor-atomization}, the same phase $\e^{\ii\Theta(t)}$ admits two exact representations.
\begin{enumerate}[label=(\alph*)]
\item The \emph{Taylor walk}
\[
\e^{\ii\Theta(t)}=\sum_{q=0}^{\infty}\frac{\bigl(\ii\Theta(t)\bigr)^q}{q!},
\]
whose step directions cycle with period $4$ and whose step lengths are $|\Theta(t)|^q/q!$.
\item The \emph{atomic Tobley walk}, obtained as the uniform limit of the normalized endpoints of the finite orthogonal walks attached to the atomic slopes $\alpha_n(t)=\tan(\beta_n t^n)$.
\end{enumerate}
\end{corollary}

\begin{proof}
Part~(a) is the Taylor expansion of the exponential function.  Part~(b) is Theorem~\ref{thm:taylor-atomization}(iii)--(iv).
\end{proof}

\section{Taylor atomization of nonvanishing holomorphic germs}

\begin{theorem}[Local phase expansion]\label{thm:nonvanishing-germ}
Let $g$ be holomorphic on a neighborhood of $0\in\CC$, and assume $g(0)\neq0$.  Then there exist $\varepsilon>0$ and a unique holomorphic function
\[
h(z)=\sum_{n=1}^{\infty} c_n z^n
\]
on $|z|<\varepsilon$ such that
\[
g(z)=g(0)\e^{h(z)},
\qquad h(0)=0.
\]
Fix a unit direction $\nu\in\Sph$, and define
\[
\Theta_{g,\nu}(t):=\Ima\,h(t\nu)
=\sum_{n=1}^{\infty}\beta_{g,\nu,n} t^n,
\qquad
\beta_{g,\nu,n}:=\Ima\bigl(c_n\nu^n\bigr).
\]
Then, for $0<t<\varepsilon$,
\[
\ph\bigl(g(t\nu)\bigr)=\ph\bigl(g(0)\bigr)\e^{\ii\Theta_{g,\nu}(t)}.
\]
Moreover, after shrinking $\varepsilon$ if necessary, Theorem~\ref{thm:taylor-atomization} applies to $\Theta_{g,\nu}$, so the local phase is represented by a canonical atomic sequence
\[
\alpha_{g,\nu,n}(t):=\tan\bigl(\beta_{g,\nu,n}t^n\bigr).
\]
\end{theorem}

\begin{proof}
Because $g(0)\neq0$, after shrinking the neighborhood we may assume that $g$ never vanishes on the disc $|z|<\varepsilon$.  Since the disc is simply connected, the nonvanishing holomorphic function $g(z)/g(0)$ admits a unique holomorphic logarithm $h$ with $h(0)=0$.  Expanding $h$ into its Taylor series gives
\[
g(z)=g(0)\e^{h(z)},
\qquad
h(z)=\sum_{n=1}^{\infty} c_n z^n.
\]
Evaluating at $z=t\nu$ yields
\[
g(t\nu)=g(0)\exp\!\Bigl(\sum_{n=1}^{\infty} c_n\nu^n t^n\Bigr).
\]
Taking moduli gives
\[
|g(t\nu)|=|g(0)|\exp\!\Bigl(\sum_{n=1}^{\infty}\Rea(c_n\nu^n)t^n\Bigr),
\]
while dividing by the modulus yields
\[
\ph\bigl(g(t\nu)\bigr)=\ph\bigl(g(0)\bigr)\exp\!\Bigl(\ii\sum_{n=1}^{\infty}\Ima(c_n\nu^n)t^n\Bigr)
=\ph\bigl(g(0)\bigr)\e^{\ii\Theta_{g,\nu}(t)}.
\]
The last claim is exactly Theorem~\ref{thm:taylor-atomization} applied to the real-analytic function $\Theta_{g,\nu}$.
\end{proof}

\begin{proposition}[The first two Taylor atoms]\label{prop:first-two-atoms}
In the situation of Theorem~\ref{thm:nonvanishing-germ}, write
\[
\frac{g(z)}{g(0)}=1+a_1 z+a_2 z^2+O(z^3)
\qquad (z\to0).
\]
Then
\[
c_1=a_1,
\qquad
c_2=a_2-\frac12 a_1^2,
\]
and therefore
\[
\beta_{g,\nu,1}=\Ima(\nu a_1),
\qquad
\beta_{g,\nu,2}=\Ima\!\left(\nu^2\Bigl(a_2-\frac12 a_1^2\Bigr)\right).
\]
\end{proposition}

\begin{proof}
Since $\log(1+w)=w-\frac12 w^2+O(w^3)$ as $w\to0$, applying this with
\[
w=a_1 z+a_2 z^2+O(z^3)
\]
gives
\[
h(z)=a_1 z+\Bigl(a_2-\frac12 a_1^2\Bigr)z^2+O(z^3).
\]
Comparing coefficients with $h(z)=\sum_{n\ge1} c_n z^n$ proves the formula for $c_1$ and $c_2$.  The formulas for the $\beta$'s are immediate from the definition $\beta_{g,\nu,n}=\Ima(c_n\nu^n)$.
\end{proof}

\section{Local phase packets of the non-trivial zeta zeros}

\begin{definition}
Let $\rho\in Z_\xi$ be a non-trivial zero of multiplicity
\[
m:=m_\rho=\ord_\rho(\xi).
\]
Define the leading coefficient
\[
\lambda_\rho:=\frac{\xi^{(m)}(\rho)}{m!}\in\CC^\times.
\]
Then the reduced local germ at $\rho$ is
\[
\psi_\rho(z):=\frac{\xi(\rho+z)}{\lambda_\rho z^m}.
\]
By construction $\psi_\rho(0)=1$.
\end{definition}

\begin{theorem}[Local phase packet at a zeta zero]\label{thm:local-phase-packet}
Let $\rho\in Z_\xi$ have multiplicity $m$, and let $\nu\in\Sph$ be a unit direction.  Then there exists $\varepsilon>0$ and a unique holomorphic function
\[
h_\rho(z)=\sum_{n=1}^{\infty} c_{\rho,n} z^n
\qquad (|z|<\varepsilon)
\]
such that
\[
\xi(\rho+z)=\lambda_\rho z^m \e^{h_\rho(z)},
\qquad h_\rho(0)=0.
\]
Define the local phase series
\[
\Theta_{\rho,\nu}(t):=\Ima\,h_\rho(t\nu)
=\sum_{n=1}^{\infty}\beta_{\rho,\nu,n} t^n,
\qquad
\beta_{\rho,\nu,n}:=\Ima\bigl(c_{\rho,n}\nu^n\bigr).
\]
Then, for $0<t<\varepsilon$,
\begin{align*}
|\xi(\rho+t\nu)|
&=t^m |\lambda_\rho| \exp\!\Bigl(\sum_{n=1}^{\infty}\Rea(c_{\rho,n}\nu^n)t^n\Bigr),\\[4pt]
\ph\bigl(\xi(\rho+t\nu)\bigr)
&=\nu^m\ph(\lambda_\rho)\e^{\ii\Theta_{\rho,\nu}(t)}.
\end{align*}
After shrinking $\varepsilon$ if necessary, the phase factor admits the exact atomic representation
\[
\e^{\ii\Theta_{\rho,\nu}(t)}
=\prod_{n=1}^{\infty}\nu\bigl(\alpha_{\rho,\nu,n}(t)\bigr),
\qquad
\alpha_{\rho,\nu,n}(t):=\tan\bigl(\beta_{\rho,\nu,n}t^n\bigr),
\]
and its finite truncations are normalized endpoints of the orthogonal walks
\[
W_{\rho,\nu,M}(t):=
\sum_{q=0}^M \ii^q e_q\bigl(\alpha_{\rho,\nu,1}(t),\dots,\alpha_{\rho,\nu,M}(t)\bigr).
\]
\end{theorem}

\begin{proof}
Because $\xi(\rho+z)$ has a zero of exact order $m$ at $z=0$, one has a factorization
\[
\xi(\rho+z)=\lambda_\rho z^m \psi_\rho(z)
\]
with $\psi_\rho$ holomorphic and $\psi_\rho(0)=1$.  Shrinking $\varepsilon$ if necessary, we may assume that $\psi_\rho$ never vanishes on $|z|<\varepsilon$.  Since the disc is simply connected, there exists a unique holomorphic logarithm $h_\rho$ with $h_\rho(0)=0$ and
\[
\psi_\rho(z)=\e^{h_\rho(z)}.
\]
This proves the factorization
\[
\xi(\rho+z)=\lambda_\rho z^m \e^{h_\rho(z)}.
\]
Now evaluate at $z=t\nu$ with $0<t<\varepsilon$.  Because $|\nu|=1$,
\[
\xi(\rho+t\nu)=\lambda_\rho t^m \nu^m \exp\!\Bigl(\sum_{n=1}^{\infty} c_{\rho,n}\nu^n t^n\Bigr).
\]
Taking moduli yields the first displayed formula, and dividing by the modulus yields
\[
\ph\bigl(\xi(\rho+t\nu)\bigr)=\nu^m\ph(\lambda_\rho)\exp\!\Bigl(\ii\sum_{n=1}^{\infty}\Ima(c_{\rho,n}\nu^n)t^n\Bigr)
=\nu^m\ph(\lambda_\rho)\e^{\ii\Theta_{\rho,\nu}(t)}.
\]
The atomic representation and the orthogonal-walk formula are exactly Theorem~\ref{thm:taylor-atomization} applied to the real-analytic function $\Theta_{\rho,\nu}$.
\end{proof}

\begin{corollary}[The first local phase atoms of a zeta zero]\label{cor:zeta-first-atoms}
Let $\rho\in Z_\xi$ have multiplicity $m$.  Then
\[
\psi_\rho(z)=1+a_{\rho,1} z+a_{\rho,2} z^2+O(z^3),
\]
where
\[
a_{\rho,1}=\frac{\xi^{(m+1)}(\rho)}{(m+1)\xi^{(m)}(\rho)},
\qquad
a_{\rho,2}=\frac{\xi^{(m+2)}(\rho)}{(m+2)(m+1)\xi^{(m)}(\rho)}.
\]
Hence, for every unit direction $\nu$,
\[
\beta_{\rho,\nu,1}=\Ima\!\left(\nu\frac{\xi^{(m+1)}(\rho)}{(m+1)\xi^{(m)}(\rho)}\right),
\]
and
\[
\beta_{\rho,\nu,2}=\Ima\!\left(\nu^2\Biggl(\frac{\xi^{(m+2)}(\rho)}{(m+2)(m+1)\xi^{(m)}(\rho)}-
\frac12\Bigl(\frac{\xi^{(m+1)}(\rho)}{(m+1)\xi^{(m)}(\rho)}\Bigr)^2\Biggr)\right).
\]
In particular, if $\rho$ is simple and lies on the critical line, then the first normal and tangent atoms are
\[
\beta_{\rho,1,1}=\Ima\!\left(\frac{\xi''(\rho)}{2\xi'(\rho)}\right),
\qquad
\beta_{\rho,\ii,1}=\Rea\!\left(\frac{\xi''(\rho)}{2\xi'(\rho)}\right).
\]
\end{corollary}

\begin{proof}
Since
\[
\xi(\rho+z)=\sum_{k=m}^{\infty}\frac{\xi^{(k)}(\rho)}{k!}z^k,
\]
dividing by $\lambda_\rho z^m$ gives
\[
\psi_\rho(z)=1+
\frac{\xi^{(m+1)}(\rho)}{(m+1)\xi^{(m)}(\rho)}z+
\frac{\xi^{(m+2)}(\rho)}{(m+2)(m+1)\xi^{(m)}(\rho)}z^2+O(z^3).
\]
Now apply Proposition~\ref{prop:first-two-atoms}.  If $m=1$ and $\nu=\ii$, then
\[
\beta_{\rho,\ii,1}=\Ima\!\left(\ii\frac{\xi''(\rho)}{2\xi'(\rho)}\right)=\Rea\!\left(\frac{\xi''(\rho)}{2\xi'(\rho)}\right).
\]
\end{proof}

\begin{proposition}[Packet symmetry of the local phase data]\label{prop:packet-symmetry-local}
Let $\rho\in Z_\xi$ have multiplicity $m$.  Then
\[
\lambda_{1-\rho}=(-1)^m\lambda_\rho,
\qquad
\lambda_{\bar\rho}=\overline{\lambda_\rho}.
\]
Moreover, after shrinking the common radius of definition if necessary, the reduced local germs satisfy
\[
\psi_{1-\rho}(z)=\psi_\rho(-z),
\qquad
\psi_{\bar\rho}(z)=\overline{\psi_\rho(\bar z)}.
\]
Consequently, the logarithmic phase series satisfy
\[
\Theta_{1-\rho,-\nu}(t)=\Theta_{\rho,\nu}(t),
\qquad
\Theta_{\bar\rho,\bar\nu}(t)=-\Theta_{\rho,\nu}(t),
\]
and therefore the atomic slopes transform by
\[
\alpha_{1-\rho,-\nu,n}(t)=\alpha_{\rho,\nu,n}(t),
\qquad
\alpha_{\bar\rho,\bar\nu,n}(t)=-\alpha_{\rho,\nu,n}(t).
\]
In particular,
\[
|\lambda_{\rho}|=|\lambda_{\bar\rho}|=|\lambda_{1-\rho}|=|\lambda_{1-\bar\rho}|,
\]
so the local jet weight
\[
\omega(P_\rho):=|\lambda_\rho|
\]
is a well-defined positive invariant of the packet $P_\rho$.
\end{proposition}

\begin{proof}
Because $\xi(1-s)=\xi(s)$, differentiating $m$ times gives
\[
(-1)^m\xi^{(m)}(1-s)=\xi^{(m)}(s).
\]
Setting $s=\rho$ yields
\[
\xi^{(m)}(1-\rho)=(-1)^m\xi^{(m)}(\rho),
\]
hence $\lambda_{1-\rho}=(-1)^m\lambda_\rho$.  Likewise, $\overline{\xi(\bar s)}=\xi(s)$ implies
\[
\xi^{(m)}(\bar\rho)=\overline{\xi^{(m)}(\rho)},
\]
whence $\lambda_{\bar\rho}=\overline{\lambda_\rho}$.

Next,
\[
\xi(1-\rho+z)=\xi(\rho-z)=\lambda_\rho(-z)^m\psi_\rho(-z)=\lambda_{1-\rho} z^m\psi_\rho(-z),
\]
so by the definition of $\psi_{1-\rho}$ one has $\psi_{1-\rho}(z)=\psi_\rho(-z)$.  Similarly,
\[
\xi(\bar\rho+z)=\overline{\xi(\rho+\bar z)}
=\overline{\lambda_\rho \bar z^m \psi_\rho(\bar z)}
=\lambda_{\bar\rho} z^m\overline{\psi_\rho(\bar z)},
\]
so $\psi_{\bar\rho}(z)=\overline{\psi_\rho(\bar z)}$.

Let $h_\rho$ be the logarithm used in Theorem~\ref{thm:local-phase-packet}.  Since $\psi_{1-\rho}(0)=1$ and $\psi_\rho(0)=1$, uniqueness of the logarithm normalized by value $0$ at the origin gives
\[
h_{1-\rho}(z)=h_\rho(-z),
\qquad
h_{\bar\rho}(z)=\overline{h_\rho(\bar z)}.
\]
Evaluating at $z=-t\nu$ and $z=t\bar\nu$ shows
\[
\Theta_{1-\rho,-\nu}(t)=\Ima\,h_{1-\rho}(-t\nu)=\Ima\,h_\rho(t\nu)=\Theta_{\rho,\nu}(t),
\]
and
\[
\Theta_{\bar\rho,\bar\nu}(t)=\Ima\,\overline{h_\rho(t\nu)}=-\Ima\,h_\rho(t\nu)=-\Theta_{\rho,\nu}(t).
\]
The slope formulas follow because $\tan$ is odd.  The weight invariance is immediate from the transformation law for the $\lambda$'s.
\end{proof}
% END part_phase

\section{Integrated conclusion}

We now record in one theorem the mathematical content of the integrated note.

\begin{theorem}[Integrated split-zero zeta theorem]\label{thm:integrated-final}
Let
\[
S=G(\ZZ)=\ZZ\sqcup\{\taup\},
\qquad
A=\{0,\taup\}\cong \BB,
\qquad
F_1(s):=\zeta(s)-1.
\]
Then the Riemann zeta function and its completed form $\xi$ carry the following exact structures.

\begin{enumerate}[label=(\roman*)]
\item For every zero $\rho$ of $\zeta$ one has $F_1(\rho)=-1$ and $\Fix_S(F_1(\rho))=A$.  Thus the reduced pointwise split-zero coefficient is constant.
\item For every effective divisor $D$ on a Riemann surface, the finite-support presheaf $\cP_D$ sheafifies to the Boolean multiplicity shadow $\cV_D$, the reduced shadow $\cA_D$ is a canonical split retract of $\cV_D$, and $\cV_D\cong \cA_D$ if and only if every multiplicity is $1$.
\item The classical local Artin thickening $Q_{x,m}=\cO_{X,x}/\mathfrak m_x^m$ has ideal lattice the terminal-segment chain $C_m$, while the Boolean multiplicity stalk is the cube $B_m\cong A^m$, and the exact comparison is the biadjunction
\[
L_m\dashv i_m\dashv R_m.
\]
\item For a holomorphic function $f$ with divisor $D$, the class of $f$ in $K_{x,m_x}=\cO_{X,x}/\mathfrak m_x^{m_x+1}$ generates the minimal nonzero ideal.  After compactification to $\PP^1(\CC)$, both the Boolean shadows and the jet shadows acquire explicit tail objects at $\infty$; for zeta these tails are nonzero.
\item The completed zeta function removes the trivial-zero sector.  Its non-trivial-zero divisor descends to packet sheaves on $Y_\xi=|D_\xi|/W$, and the trivial zeros carry the exact weights
\[
|\zeta'(-2n)|=\frac{(2n)!}{2(2\pi)^{2n}}\,\zeta(2n+1).
\]
\item The centered completed zeta function $\Xi(z)=\xi(\tfrac12+z)$ admits an exact normalized packet product over $Y_\xi$.  Equivalently,
\[
\xi(s)=\xi\!\left(\frac12\right)\prod_{O\in Y_\xi}P_O(s)^{m_O},
\]
where $P_O$ is quadratic exactly for critical-line packets and quartic exactly for off-line packets.  Hence RH is equivalent to the absence of quartic packet factors.
\item The same packet product combines with the reciprocal gamma product to yield the global pole-trivial-packet factorization of $\zeta$, and logarithmic differentiation gives the packet-grouped partial-fraction expansion of $\zeta'/\zeta$.
\item For every non-trivial zero $\rho$ of multiplicity $m$ and every unit direction $\nu\in\Sph$, the reduced local germ of $\xi$ has a canonical logarithmic phase series
\[
\Theta_{\rho,\nu}(t)=\sum_{n\ge 1}\beta_{\rho,\nu,n}t^n,
\]
and an exact Taylor-atomic factorization
\[
\e^{\ii\Theta_{\rho,\nu}(t)}=\prod_{n\ge 1}\frac{1+\ii\tan(\beta_{\rho,\nu,n}t^n)}{\sqrt{1+\tan^2(\beta_{\rho,\nu,n}t^n)}}.
\]
Its finite truncations are normalized endpoints of explicit orthogonal walks governed by elementary symmetric polynomials of the atomic slopes.
\item In the Connes--Consani spectral coordinate $\Phi(s)=i(s-\tfrac12)$, every critical packet determines the primitive $\zeta$-cycle length $2\pi/|\gamma|$ for $\rho=\tfrac12+i\gamma$.
\end{enumerate}

Accordingly, a non-trivial zeta zero is represented in the integrated framework by a packet point in the Connes spectral plane, a Boolean multiplicity cube, a packet jet weight, and a canonical Taylor-atomic phase geometry.
\end{theorem}

\begin{proof}
Items (i)--(iv) are proved in Sections~2--6.  Item (v) is Section~6 together with Proposition~\ref{prop:trivial-derivative-and-Lambda}.  Items (vi)--(vii) are Sections~8--13.  Item (viii) is Sections~16--18.  Item (ix) is Sections~14--15.
\end{proof}

\begin{remark}
The only genuinely separate strand not built directly into Theorem~\ref{thm:integrated-final} is the spectral-regularization correction recorded in Appendix~\ref{app:wick}.  It does not alter the divisor or packet geometry of zeta zeros, but it provides a corrected spectral-zeta companion formalism and a direct odd-theta/odd-square-zeta bridge on the torus.
\end{remark}

\appendix

\section{Corrected Wick rotation, twisted spectral zeta functions, and the torus odd-theta bridge}\label{app:wick}

This appendix integrates the mathematically correct part of the uploaded modular-Hamiltonian note together with the pasted correction.  The output needed for the present consolidated note is threefold.

\begin{enumerate}[label=(\arabic*)]
\item The vacuum coefficient of a Wick-rotated trace is the ground-state coefficient of the normalized $q$-character, not the value at $s=0$ of the twisted spectral zeta function.
\item The naive identity $\mathcal L_U(0)=-c$ is false in general.
\item On the circle, the heat evolution of the basic half-period odd step function gives a direct bridge between an odd theta character and the Dirichlet series $2(1-2^{-2s})\zeta(2s)$.
\end{enumerate}

\subsection{Ground-state normalized Wick characters}

\begin{theorem}[Ground-state normalized Wick character]\label{thm:wick-character}
Let $L$ be a self-adjoint elliptic operator on a closed manifold with compact resolvent and lower-bounded spectrum
\[
\lambda_{\min}=\inf\operatorname{Spec}(L)>-\infty.
\]
Let $U$ be unitary and commute with $L$, and write $H_a:=aL$ with $a>0$.  For each eigenspace $E_\lambda$ set
\[
a_\lambda:=\operatorname{Tr}(U|_{E_\lambda}).
\]
Then for every $\tau\in\mathfrak H:=\{\tau\in\CC:\Im(\tau)>0\}$ the series
\[
\chi_{U,a}(\tau):=\operatorname{Tr}(Ue^{2\pi i\tau H_a})
\]
converges absolutely and defines a holomorphic function on $\mathfrak H$.  Moreover,
\[
e^{-2\pi i\tau a\lambda_{\min}}\chi_{U,a}(\tau)
=
\sum_{\lambda\in\operatorname{Spec}(L)} a_\lambda q^{a(\lambda-\lambda_{\min})},
\qquad q=e^{2\pi i\tau}.
\]
If the ground eigenspace is one-dimensional and $U\psi_0=c\psi_0$, then
\[
e^{-2\pi i\tau a\lambda_{\min}}\chi_{U,a}(\tau)
=
c+\sum_{\lambda>\lambda_{\min}} a_\lambda q^{a(\lambda-\lambda_{\min})}.
\]
\end{theorem}

\begin{proof}
Because $UL=LU$, the operator $Ue^{2\pi i\tau aL}$ is diagonal in a spectral basis.  Hence
\[
\chi_{U,a}(\tau)=\sum_{\lambda} a_\lambda e^{2\pi i\tau a\lambda}.
\]
Now
\[
\bigl|e^{2\pi i\tau a\lambda}\bigr|=e^{-2\pi(\Im\tau)a\lambda}
=e^{-2\pi(\Im\tau)a\lambda_{\min}}e^{-2\pi(\Im\tau)a(\lambda-\lambda_{\min})}.
\]
Therefore
\[
\sum_\lambda |a_\lambda|\,|e^{2\pi i\tau a\lambda}|
\le e^{-2\pi(\Im\tau)a\lambda_{\min}}
\sum_\lambda \dim(E_\lambda)e^{-2\pi(\Im\tau)a(\lambda-\lambda_{\min})}.
\]
The last sum is the heat trace of the nonnegative compact-resolvent operator $L-\lambda_{\min}I$ at time $2\pi a\Im\tau$, hence is finite.  Thus the series converges absolutely and locally uniformly on $\mathfrak H$.  Factoring out the ground energy gives the displayed expansion, and the one-dimensional ground-state formula is immediate.
\end{proof}

\begin{corollary}\label{cor:genuine-q-character}
Assume that there exists $a>0$ such that $a(\lambda-\lambda_{\min})\in\ZZ_{\ge 0}$ for every eigenvalue $\lambda$ of $L$.  Then the normalized trace of Theorem~\ref{thm:wick-character} is a genuine $q$-character whose vacuum coefficient is exactly the ground-state character $c$.
\end{corollary}

\begin{proof}
Group together equal integral exponents in the expansion of Theorem~\ref{thm:wick-character}.
\end{proof}

\subsection{Twisted spectral Dirichlet series and the failure of $\mathcal L_U(0)=-c$}

\begin{theorem}[Correct Mellin formula for the twisted spectral Dirichlet series]\label{thm:twisted-mellin}
Assume $L\ge 0$, $\dim\ker L=1$, and $U\psi_0=c\psi_0$.  Define
\[
\Theta_U(t):=\operatorname{Tr}(Ue^{-tL}),
\qquad
\mathcal L_U(s):=\operatorname{Tr}'(UL^{-s})=\sum_{\lambda>0} a_\lambda\lambda^{-s},
\]
where the prime indicates omission of the kernel.  Then for $\Re s$ sufficiently large,
\[
\mathcal L_U(s)=\frac{1}{\Gamma(s)}\int_0^\infty t^{s-1}\bigl(\Theta_U(t)-c\bigr)\,dt.
\]
\end{theorem}

\begin{proof}
By spectral decomposition,
\[
\Theta_U(t)=c+\sum_{\lambda>0} a_\lambda e^{-t\lambda}.
\]
For $\Re s\gg 0$, absolute convergence allows termwise integration, and
\[
\int_0^\infty t^{s-1}e^{-t\lambda}\,dt=\Gamma(s)\lambda^{-s}.
\]
Therefore
\[
\frac{1}{\Gamma(s)}\int_0^\infty t^{s-1}(\Theta_U(t)-c)\,dt
=\sum_{\lambda>0} a_\lambda\lambda^{-s}
=\mathcal L_U(s).
\]
\end{proof}

\begin{proposition}[Explicit counterexample]\label{prop:counterexample-L0}
Let $L=-\partial_x^2$ on $\mathbb T=\RR/\ZZ$, with Fourier basis $e_k(x)=e^{2\pi ikx}$, and define $U$ by
\[
Ue_0=e_0,
\qquad
Ue_k=-e_k\quad (k\ne 0).
\]
Then $U$ is unitary, $UL=LU$, the kernel character is $c=1$, but
\[
\mathcal L_U(0)=1\neq -1=-c.
\]
\end{proposition}

\begin{proof}
The eigenvalues are $Le_k=4\pi^2k^2e_k$.  Hence
\[
\mathcal L_U(s)=\sum_{k\ne 0}(-1)(4\pi^2k^2)^{-s}=-2(4\pi^2)^{-s}\zeta(2s).
\]
Evaluating at $s=0$ gives
\[
\mathcal L_U(0)=-2\zeta(0)=-2\Bigl(-\frac12\Bigr)=1.
\]
\end{proof}

\begin{remark}
Theorem~\ref{thm:wick-character} and Proposition~\ref{prop:counterexample-L0} isolate the correction exactly: the vacuum coefficient is read off from the ground-state normalized Wick character, while the value at $s=0$ of the twisted spectral zeta function is controlled by short-time heat coefficients.
\end{remark}

\subsection{The torus spline tower and the odd theta sector}

Let
\[
h_0(x):=\mathbf 1_{[0,1/2)}(x)-\mathbf 1_{[1/2,1)}(x)
\qquad (x\in\mathbb T).
\]
For a mean-zero function $f\in L^2_0(\mathbb T)$ define the mean-zero antiderivative $Q=D^{-1}f$ by Fourier coefficients
\[
\widehat{Qf}(0)=0,
\qquad
\widehat{Qf}(k)=\frac{\widehat f(k)}{2\pi i k}\quad (k\ne 0).
\]
Set
\[
g_0:=h_0,
\qquad
g_n:=Q^n g_0\quad (n\ge 1).
\]

\begin{theorem}[Exact Fourier formula for the spline tower]\label{thm:spline-fourier}
For every $n\ge 0$ and every $k\ne 0$ one has
\[
\widehat g_n(k)=2\,\frac{1-(-1)^k}{(2\pi i k)^{n+1}}.
\]
Equivalently,
\[
\widehat g_n(2m)=0,
\qquad
\widehat g_n(2m+1)=\frac{4}{(2\pi i(2m+1))^{n+1}}.
\]
Moreover,
\[
g_n\Bigl(x+\frac12\Bigr)=-g_n(x)
\qquad (n\ge 0),
\]
so the single chain $(g_n)_{n\ge 0}$ lies in the proper closed half-period odd subspace
\[
L^2_-(\mathbb T):=\bigl\{f\in L^2(\mathbb T):f(x+1/2)=-f(x)\bigr\}.
\]
In particular, $(g_n)_{n\ge 0}$ is not a basis of all $L^2(\mathbb T)$.
\end{theorem}

\begin{proof}
For $n=0$, direct computation gives
\[
\widehat g_0(k)=\int_0^{1/2}e^{-2\pi ikx}\,dx-\int_{1/2}^{1}e^{-2\pi ikx}\,dx
=\frac{1-(-1)^k}{\pi i k}
=2\,\frac{1-(-1)^k}{2\pi i k}.
\]
If the formula holds for $n$, then since $g_{n+1}=Qg_n$,
\[
\widehat g_{n+1}(k)=\frac{\widehat g_n(k)}{2\pi i k}
=2\,\frac{1-(-1)^k}{(2\pi i k)^{n+2}},
\]
which proves the first claim by induction.  Because only odd frequencies occur, translating by $1/2$ multiplies every Fourier mode by $-1$, hence $g_n(x+1/2)=-g_n(x)$.  The closed span of the $g_n$ is therefore contained in the proper closed subspace $L^2_-(\mathbb T)$, so it cannot equal all of $L^2(\mathbb T)$.
\end{proof}

\begin{theorem}[Distributional spline equation and Sobolev threshold]\label{thm:spline-regularity}
For every $n\ge 0$ one has
\[
D^{n+1}g_n=2(\delta_0-\delta_{1/2})
\]
in the sense of periodic distributions.  Moreover,
\[
g_n\in H^s(\mathbb T)
\quad\Longleftrightarrow\quad
s<n+\frac12.
\]
\end{theorem}

\begin{proof}
For $n=0$, differentiating the jump step gives $Dg_0=2(\delta_0-\delta_{1/2})$.  Since $Dg_n=g_{n-1}$ for $n\ge 1$, repeated differentiation yields the stated distributional identity.

For the Sobolev statement, Theorem~\ref{thm:spline-fourier} gives $|\widehat g_n(k)|\asymp |k|^{-(n+1)}$ on odd frequencies and zero on even frequencies.  Hence
\[
\|g_n\|_{H^s}^2\asymp \sum_{m\in\ZZ\setminus\{0\}} (1+|2m+1|^2)^s |2m+1|^{-2n-2}.
\]
This converges exactly when $2s-2n-2<-1$, i.e. when $s<n+\tfrac12$.
\end{proof}

\begin{theorem}[Integrated Haar basis theorem]\label{thm:integrated-haar}
Let $(\psi_{j,m})_{j\ge 0,\ 0\le m<2^j}$ be the periodized Haar orthonormal basis of $L^2_0(\mathbb T)$, with $\psi_{0,0}=h_0$.  Then for each $n\ge 0$ the family
\[
\Psi^{(n)}_{j,m}:=Q^n\psi_{j,m}
\]
is a Riesz basis of
\[
H^n_0(\mathbb T):=\{f\in H^n(\mathbb T):\widehat f(0)=0\}.
\]
Each $\Psi^{(n)}_{j,m}$ is a periodic spline of degree $n$ with dyadic knots, and the chain $(g_n)$ is exactly the scale-$0$ branch $\Psi^{(n)}_{0,0}$.
\end{theorem}

\begin{proof}
The operator $Q^n:L^2_0(\mathbb T)\to H^n_0(\mathbb T)$ is a bounded isomorphism because on nonzero Fourier modes it multiplies by $(2\pi i k)^{-n}$, and the Sobolev norm on $H^n_0$ is equivalent to the $L^2$ norm after this multiplier.  The image of an orthonormal basis under a bounded invertible operator is a Riesz basis.  Since the Haar wavelets are piecewise constant with dyadic jumps, repeated mean-zero integration produces periodic splines of degree $n$ with the same dyadic knot set.  The identity $\psi_{0,0}=h_0$ gives $\Psi^{(n)}_{0,0}=g_n$.
\end{proof}

\subsection{Odd theta characters and the Dirichlet series $2(1-2^{-2s})\zeta(2s)$}

Rescale the circle Laplacian to
\[
N:=-(4\pi^2)^{-1}\partial_x^2,
\qquad
Ne_k=k^2e_k,
\qquad e_k(x)=e^{2\pi ikx}.
\]
Let $T$ be translation by one half-period,
\[
(Tf)(x):=f\!\left(x+\frac12\right),
\]
so that $Te_k=(-1)^k e_k$, and let $P_-:=\frac12(I-T)$.

\begin{proposition}[Odd theta character]\label{prop:odd-theta-character}
For $q=e^{2\pi i\tau}$ with $\Im\tau>0$ one has
\[
\chi_-(\tau):=\operatorname{Tr}(P_-e^{2\pi i\tau N})
=\sum_{k\ \mathrm{odd}} q^{k^2}
=2\sum_{m\ge 0} q^{(2m+1)^2}.
\]
\end{proposition}

\begin{proof}
In the Fourier basis, $e^{2\pi i\tau N}e_k=q^{k^2}e_k$ and $P_-e_k=\frac{1-(-1)^k}{2}e_k$.  Therefore
\[
\operatorname{Tr}(P_-e^{2\pi i\tau N})
=\sum_{k\in\ZZ} \frac{1-(-1)^k}{2} q^{k^2}
=\sum_{k\ \mathrm{odd}} q^{k^2}.
\]
Pairing $k$ and $-k$ gives the second expression.
\end{proof}

\begin{theorem}[Spline field as a spatial primitive of the odd theta character]\label{thm:spline-theta}
Define
\[
\Phi_0(x,\tau):=e^{2\pi i\tau N}h_0(x).
\]
Then
\[
\partial_x\Phi_0(x,\tau)=4\sum_{k\ \mathrm{odd}} q^{k^2}e^{2\pi ikx},
\]
and in particular
\[
\chi_-(\tau)=\frac14\,\partial_x\Phi_0(0,\tau).
\]
\end{theorem}

\begin{proof}
By Theorem~\ref{thm:spline-fourier},
\[
\Phi_0(x,\tau)=\sum_{k\ \mathrm{odd}} \frac{2}{\pi i k}q^{k^2}e^{2\pi ikx}.
\]
Differentiate termwise to obtain
\[
\partial_x\Phi_0(x,\tau)=4\sum_{k\ \mathrm{odd}} q^{k^2}e^{2\pi ikx}.
\]
Setting $x=0$ and using Proposition~\ref{prop:odd-theta-character} yields the identity.
\end{proof}

\begin{proposition}[Odd-square Dirichlet series from the spline field]\label{prop:odd-square-dirichlet}
Let
\[
u_0(x,t):=e^{-tN}h_0(x)
=\sum_{k\ \mathrm{odd}}\frac{2}{\pi i k}e^{-tk^2}e^{2\pi ikx},
\qquad t>0.
\]
Then for $\Re s>\tfrac12$ one has
\[
\frac{1}{4\Gamma(s)}\int_0^\infty t^{s-1}\partial_x \nu_0(0,t)\,dt
=2(1-2^{-2s})\zeta(2s).
\]
\end{proposition}

\begin{proof}
From the Fourier series,
\[
\partial_x\nu_0(0,t)=4\sum_{k\ \mathrm{odd}} e^{-tk^2}.
\]
Therefore
\[
\frac{1}{4\Gamma(s)}\int_0^\infty t^{s-1}\partial_x \nu_0(0,t)\,dt
=\frac{1}{\Gamma(s)}\sum_{k\ \mathrm{odd}}\int_0^\infty t^{s-1}e^{-tk^2}\,dt
=\sum_{k\ \mathrm{odd}}k^{-2s}.
\]
Finally,
\[
\sum_{k\ \mathrm{odd}}k^{-2s}=2\sum_{m\ge 0}(2m+1)^{-2s}=2(1-2^{-2s})\zeta(2s).
\]
\end{proof}

\begin{remark}
Theorem~\ref{thm:spline-theta} and Proposition~\ref{prop:odd-square-dirichlet} are the exact torus realization hidden in the corrected note: the basic half-period odd spline field is not itself the theta character but a canonical spatial primitive of it, and the Mellin transform of its derivative produces the odd-square Dirichlet series.  This is the clean point at which the corrected spectral appendix meets the zeta side of the present note.
\end{remark}

\begin{thebibliography}{99}

\bibitem{ConnesConsaniCycles}
A.~Connes and C.~Consani,
\emph{Hochschild homology, trace map and $\zeta$-cycles},
in \emph{Cyclic Cohomology at 40: Achievements and Future Prospects},
Proc. Sympos. Pure Math. \textbf{105}, American Mathematical Society, Providence, RI, 2023, pp.~83--110.

\bibitem{Lee}
J.~M.~Lee,
\emph{Introduction to Complex Manifolds},
Graduate Studies in Mathematics \textbf{244}, American Mathematical Society, Providence, RI, 2024.

\bibitem{Lorscheid}
O.~Lorscheid,
\emph{Blueprints---towards absolute arithmetic},
J. Number Theory \textbf{144} (2014), 408--421.

\bibitem{Tobley}
J.~Tobley,
\emph{Unpublished correspondence on orthogonal-walk phase decompositions, Taylor-series atomization, and affine cyclotomic charts},
project source transcripts communicated with the manuscript materials, May 2026.

\end{thebibliography}

\end{document}
